\section{Number-theory consequences with the hypotheses separated}
\label{sec:number-theory}

Throughout this section $A$ denotes the user-designated Alp\"oge theorem
for the Riemann zeta function and Dirichlet L-functions.  It is an adopted
cited input, not a claim that this run rebuilt the unseen Lean dependencies.
$M$ denotes the conditional all-number-field unitary Hecke-GRH assertion
of Masullo, \S1.2.  Every use of $M$ below remains conditional.  The
accompanying catalogue is \texttt{work/number\_theory/rows.json}; its
literature records include retrieval evidence and distinguish a checked
statement from a reread of an entire classical proof.

There are three different geometric translations.  Rational primes are
the closed points of $\operatorname{Spec}\mathbb Z$; prime ideals are the
closed points of $\operatorname{Spec}\mathcal O_K$; and a conjugacy condition
on Frobenius is a condition on a finite \'{e}tale cover of an open arithmetic
curve.  The last dictionary is the Galois-category equivalence in SGA~1,
Expos\'e~V, \S7.  The class group is
$\operatorname{Pic}(\operatorname{Spec}\mathcal O_K)$ and its class fields
are finite \'{e}tale covers.  These translations identify the objects to
which the estimates apply.  They do not turn number-field GRH into a
consequence of finite-field purity.  The inherited cote-1 trace programme
is a comparison of explicit formulas; no new archive scan is certified by
this section.  Exact archive work belongs to the geometry section.

\subsection{The classical arithmetic bounds being used}

\paragraph{NT-P01--NT-P04: rational primes and gaps.}
Schoenfeld, \emph{Sharper bounds for the Chebyshev functions
$\theta(x)$ and $\psi(x)$. II} (1976), Theorem~10, printed p.~337,
equations (6.2) and (6.3), gives, under RH,
\[
 |\psi(x)-x|<\frac{\sqrt{x}\log^2x}{8\pi}\quad(x\ge73.2),\qquad
 |\theta(x)-x|<\frac{\sqrt{x}\log^2x}{8\pi}\quad(x\ge599).
\]
Equation (6.1) replaces $\log^2x$ by $(\log x-2)\log x$ for
$x\ge23\cdot10^8$.  Corollary~1 gives
$|\pi(x)-\operatorname{li}(x)|<\sqrt{x}\log x/(8\pi)$ for
$x\ge2657$.  Here $\operatorname{li}$ has its usual normalization.
The constants in Theorem~10 were checked in the primary indexed text.
An independently retrieved primary proof locator for the same
prime-counting bound is Johnston (2022), Theorem~3.2, equation (3.6),
\emph{The Ramanujan Journal} \textbf{59}, 1307--1321.
It requires RH only to a height $T$ satisfying
$9.06\sqrt{x/\log x}/\log\log x\le T$; under $A$ one may choose
such a $T$ for each $x$.  Its \S2 explicitly says its main theorem applies
also to the ordinary, non-half-weight counting functions.  Equation (3.3)
there also improves the threshold for the displayed $\psi$ inequality
to $x\ge59$.  Thus the inaccessible original Corollary~1 page is no
longer a proof-source gate for NT-P03.

Dudek--Greni\'e--Molteni, \emph{Primes in explicit short intervals on RH}
(2016), Theorem~1.1, author PDF p.~1, gives a prime in
\[
 (x-c\sqrt{x}\log x,\ x+c\sqrt{x}\log x),
 \qquad c=\tfrac12+2/\log x,\qquad x\ge2.
\]
The same theorem gives at least $\sqrt{x}$ primes after replacing $c$ by
$c+1$.  These results imply the classical
$p_{n+1}-p_n=O(\sqrt{p_n}\log p_n)$ consequence of RH.  Neither this
bound nor the assertion of RH supplies a proof here of Cram\'er's
$O(\log^2p_n)$ conjecture, twin primes, or the binary Goldbach conjecture.
Those problems require additional information about prime correlations.

\paragraph{NT-P05--NT-P08: progressions and nonresidues.}
Greni\'e--Molteni--Perelli (2016), Theorem~1, author PDF p.~2, proves that
absolute $x_0,c_1,c_2>0$ work in
\[
 \pi(x+h;q,a)-\pi(x;q,a)\ge
 c_2\frac{h}{\varphi(q)\log x},\qquad
 c_1\varphi(q)\sqrt{x}\log x\le h\le x,\quad x\ge x_0,
\]
assuming the Dirichlet GRH modulo $q$ and $(a,q)=1$.  The following
sentence in that source gives the asymptotic if the ratio to the lower
scale tends to infinity.  We have not converted the unspecified absolute
constants into numerical constants.

Lamzouri--Li--Soundararajan (LLS), corrected
\texttt{arXiv:1309.3595v3}, is used for the sharp explicit conductor
statements.  Corollary~1.2, author p.~3, gives the least prime in a reduced
class modulo $q>3$ at most $(\varphi(q)\log q)^2$.
Corollary~1.1, p.~2, gives a quadratic nonresidue below $(\log q)^2$ for
prime $q\ge5$.  Theorem~1.1(1), p.~2, gives the proper-subgroup bound
recorded with its explicit functions $A(q),B(q)$ in NT-P08.
These are the already carried IA02 consequences, with the corrected
version and hypotheses rechecked; they are not new findings of this run.
Geometrically the progression is a Frobenius class in a cyclotomic cover,
and the nonresidue is detected by a finite rank-one character.

\paragraph{NT-P09: the least primitive root is a different quantifier.}
Shoup, \emph{Searching for primitive roots in finite fields} (1992),
Theorem~3, author PDF p.~4, gives under Dirichlet GRH
\[
 g(p)\ll r^4(\log r+1)^4(\log p)^2,
 \qquad r=\omega(p-1).
\]
The printed theorem was checked on the page image.  Here is the short
simplification to $g(p)\ll\log^6p$.  If $q_1<\cdots<q_r$ are the distinct
prime divisors of $p-1$, then $q_j\ge j+1$ and therefore
$p-1\ge\prod_{j=1}^r(j+1)$.  For $r\ge4$ the last at least $r/2$ factors
are at least $r/2$, so
$\log(p-1)\ge(r/2)\log(r/2)$.  Consequently
$r(\log r+1)\ll\log p$; enlarging the absolute constant handles $r<4$.
Substitution proves the assertion.  This finds a small generator for a
fixed $\mathbb F_p^\times$.  The density theorem for one fixed integer
as $p$ varies uses different, generally nonabelian, fields (NT-A01).

\subsection{A quantitative Frobenius trace transfer}

\paragraph{NT-C01 and NT-C02: the precise cited Chebotarev input.}
Let $L/K$ be a finite Galois extension, $G=\operatorname{Gal}(L/K)$,
$n_L=[L:\mathbb Q]$, and $\Delta_L=|\operatorname{disc}L|$.
For a conjugacy class $C$ put $\delta_C=|C|/|G|$.  For unramified
primes only, define
\[
 \psi_C(x)=\sum_{\substack{\mathfrak p\text{ unramified},\ m\ge1\\
                    (N\mathfrak p)^m\le x}}
 \mathbf1_C(\operatorname{Frob}_{\mathfrak p}^{m})\log N\mathfrak p,
 \qquad
 \pi_C(x)=\sum_{\substack{\mathfrak p\text{ unramified}\\N\mathfrak p\le x}}
             \mathbf1_C(\operatorname{Frob}_{\mathfrak p}).
\]
Set
\begin{align*}
 E_L(x)&=\sqrt{x}\left[
       \left(\frac{\log x}{2\pi}+2\right)\log\Delta_L+
       \left(\frac{\log^2x}{8\pi}+2\right)n_L\right],\qquad x\ge1,\\
 B_L(x)&=\sqrt{x}\left[
       \left(\frac1{2\pi}+\frac3{\log x}\right)\log\Delta_L+
       \left(\frac{\log x}{8\pi}+\frac1{4\pi}+\frac6{\log x}\right)n_L
                         \right],\qquad x\ge2,\\
 \operatorname{Li}_2(x)&=\int_2^x\frac{dt}{\log t}.
\end{align*}
Greni\'e--Molteni, \emph{An explicit Chebotarev density theorem under GRH}
(2019), Theorem~1.1 and Corollary~1.2, author PDF pp.~2--3, give
\[
 |\psi_C(x)-\delta_Cx|\le\delta_C E_L(x),\qquad
 |\pi_C(x)-\delta_C\operatorname{Li}_2(x)|\le\delta_C B_L(x),
\tag{NT-C}
\]
under RH for $\zeta_L$.  The theorem explicitly includes the unramified
version of $\psi_C$.  Its constants and ranges were checked on the page
images.  In particular Artin holomorphy is not an additional hypothesis.
$M$ supplies this input for every $L$; $A$ supplies it if $L/\mathbb Q$
is abelian by the factorization into primitive Dirichlet functions.

\paragraph{NT-F01. Statement and full proof.}
For a complex class function $f$ on $G$ put
\[
 \mu_f=\frac1{|G|}\sum_{g\in G}f(g),\qquad
 \|f\|_{1,G}=\frac1{|G|}\sum_{g\in G}|f(g)|,
\]
and define $\Psi_f$ and $\Pi_f$ by replacing $\mathbf1_C$ in the two
sums above by $f$.  Then
\[
 |\Psi_f(x)-\mu_fx|\le\|f\|_{1,G}E_L(x),\qquad
 |\Pi_f(x)-\mu_f\operatorname{Li}_2(x)|\le\|f\|_{1,G}B_L(x).
\tag{NT-F01}
\]
If $f=\chi_V$ is the character of a finite-dimensional complex
representation, then $\mu_f=\dim V^G$ and
$\|f\|_{1,G}\le\sqrt{\langle\chi_V,\chi_V\rangle_G}$.

To prove the estimates, choose one representative $g_C$ from each
conjugacy class.  The finite identities
\[
 f=\sum_C f(g_C)\mathbf1_C,
 \qquad \mu_f=\sum_C\delta_Cf(g_C)
\]
give $\Psi_f-\mu_fx=\sum_Cf(g_C)(\psi_C-\delta_Cx)$.
The triangle inequality and (NT-C) bound its modulus by
$E_L(x)\sum_C\delta_C|f(g_C)|=E_L(x)\|f\|_{1,G}$.
The same finite argument proves the assertion for $\Pi_f$.
The averaging operator $|G|^{-1}\sum_g\rho(g)$ is the projection to $V^G$,
so its trace is $\mu_{\chi_V}=\dim V^G$.  Cauchy--Schwarz for the
probability measure assigning mass $1/|G|$ to each group element yields
\[
 \left(\frac1{|G|}\sum_g|\chi_V(g)|\right)^2
 \le\frac1{|G|}\sum_g|\chi_V(g)|^2
 =\langle\chi_V,\chi_V\rangle_G.
\]
This proves every assertion.  This is a linear functional estimate on the
complexified representation ring; virtual differences need not be
nonnegative, so their error is controlled by the displayed norm.

\paragraph{NT-F02. Finite \texorpdfstring{\'{e}}{e}tale point counts, with the boundary.}
Let $U$ be a nonempty open subscheme of
$\operatorname{Spec}\mathcal O_K$, and let $X\to U$ be finite \'{e}tale
of degree $d$.  Let $S_X$ be the set of geometric points of its generic
fibre, let $L/K$ be its finite Galois splitting field, and let $r$ be the
number of $G$-orbits on $S_X$, equivalently the number of generic connected
components.  Let $S$ be the finite set of primes not in $U$ which are
unramified in $L$.  Under RH for $\zeta_L$,
\[
 \left|\sum_{\substack{\mathfrak p\in U,\ m\ge1\\(N\mathfrak p)^m\le x}}
       \#X(k_{\mathfrak p^m})\log N\mathfrak p-rx\right|
 \le rE_L(x)+d|S|\log x,\qquad x\ge1.
\tag{NT-F02}
\]
Here $k_{\mathfrak p^m}$ denotes the degree-$m$ extension of
$k_{\mathfrak p}$.  In particular, when no additional unramified prime is
removed, the error is $rE_L(x)$, independent of the degree $d$.

For the proof, finite \'{e}taleness implies that inertia at a prime of $U$
acts trivially on the generic geometric fibre.  The action of
$\operatorname{Gal}(L/K)$ on this fibre is faithful by the definition of
the splitting field, so $L/K$ is unramified at every prime of $U$.
Let $\chi(g)=\#S_X^g$ be the permutation character.  The elementary
Frobenius fixed-point formula for a finite \'{e}tale scheme over a finite
field is
$\#X(k_{\mathfrak p^m})=\chi(\operatorname{Frob}_{\mathfrak p}^{m})$:
its geometric points rational over that field are precisely the points
fixed by the $m$th Frobenius power.  Also $0\le\chi(g)\le d$.
Double counting pairs $(g,s)$ with $gs=s$ gives
\[
 \sum_{g\in G}\chi(g)=\sum_{s\in S_X}|G_s|=r|G|.
\]
Thus both $\mu_\chi$ and $\|\chi\|_{1,G}$ equal $r$.  NT-F01 bounds
the sum over every unramified prime by $rx$ with error $rE_L(x)$.
For each omitted prime, the nonnegative lost contribution is at most
\[
 d\left\lfloor\frac{\log x}{\log N\mathfrak p}\right\rfloor
          \log N\mathfrak p\le d\log x.
\]
If $N\mathfrak p>x$ the contribution is zero.  Summing this bound over
$S$ and applying the triangle inequality proves (NT-F02), including
$x=1$.  Notice that replacing the nonnegative character norm by the square
root of its second moment would give a weaker bound; the orbit count is
the relevant parameter here.

\paragraph{NT-F03. Splitting types.}
For a partition $\lambda$ of $d=|S_X|$, let $D_\lambda\subseteq G$ be
the elements with cycle partition $\lambda$ on $S_X$, and put
$\delta_\lambda=|D_\lambda|/|G|$.  This is a union of conjugacy classes.
Sum the second inequality (NT-C) over these classes and use the triangle
inequality.  Their densities sum to $\delta_\lambda$, proving
\[
 \left|\#\{\mathfrak p:\ N\mathfrak p\le x,\ \mathfrak p
 \text{ unramified},\ \operatorname{Frob}_{\mathfrak p}\in D_\lambda\}
 -\delta_\lambda\operatorname{Li}_2(x)\right|
 \le\delta_\lambda B_L(x).
\]
If $D_\lambda$ is empty both sides before the absolute value are zero.
For the canonical finite \'{e}tale model of the generic algebra over the
unramified locus, each Frobenius orbit is one closed point of its fibre,
and its orbit length is that point's residue degree.  Hence this cycle
partition is exactly its splitting type.  For the originally specified
model $X$ over $U$, restrict the counted primes to $U$; at most $|S|$
terms have been removed, so the same estimate holds with the additional
error $|S|$.  This proves both interpretations without assigning a fibre
of $X$ outside its specified base.

\paragraph{NT-F04. Avoiding a prescribed finite set.}
Let $D\subseteq G$ be a nonempty conjugacy-stable set of density
$\delta_D$, and let $S$ be any finite set of primes of $K$.  The same
finite-union argument gives
$\pi_D(x)\ge\delta_D(\operatorname{Li}_2(x)-B_L(x))$.
At most $|S|$ of the counted primes belong to $S$.  If
\[
 \delta_D\bigl(\operatorname{Li}_2(x)-B_L(x)\bigr)>|S|,
\]
the count outside $S$ is strictly positive.  It is an integer, so an
eligible prime exists.  This proves the avoidance assertion without adding
an unsupported arbitrary-exclusion clause to Bach--Sorenson's least-prime
theorem.  That theorem is NT-C03, Theorem~5.1, printed p.~1729, with
RH hypotheses specified in its \S2, p.~1720.

\subsection{Class numbers, Kummer towers, and Euclidean algorithms}

\paragraph{NT-H01--NT-H04: the regulator normalization.}
Write
\begin{align*}
 U(q)&=2e^\gamma\left(\log\log q-\log2+\tfrac12+
                                      \frac1{\log\log q}\right),\\
 V(q)&=\frac{12e^\gamma}{\pi^2}\left(\log\log q-\log2+\tfrac12+
                   \frac1{\log\log q}+\frac{14\log\log q}{\log q}\right).
\end{align*}
LLS Theorem~1.5, author pp.~3--4, gives
$V(q)^{-1}\le|L(1,\chi)|\le U(q)$ for primitive conductor
$q\ge10^{10}$.  Its Corollary~1.3 is the already carried imaginary
quadratic bound in NT-H02.  The inequality is not a two-sided estimate of
order $\log\log q$: the reciprocal lower bound must be retained.

Here is the full deduction of NT-H03, including its normalization.
For $D>0$ fundamental put $K=\mathbb Q(\sqrt D)$ and let $R_D$ be the
regulator of the full unit group modulo its torsion, so
$R_D=\log\varepsilon$ for a fundamental unit $\varepsilon>1$, regardless
of its norm's sign.  The analytic class-number formula is
\[
 \mathop{\rm Res}_{s=1}\zeta_K(s)=\frac{2h_DR_D}{\sqrt D}.
\]
An exact independent locator is Milne, \emph{Class Field Theory}, v4.03,
Chapter~VI, Corollary~2.12 and Example~2.13, printed p.~191.
Since $\zeta_K(s)=\zeta(s)L(s,\chi_D)$ and $\zeta$ has residue one,
$L(1,\chi_D)=2h_DR_D/\sqrt D>0$.
Multiplying the two inequalities in NT-H01 by $\sqrt D/2$ proves
\[
 \frac{\sqrt D}{2V(D)}\le h_DR_D\le\frac{\sqrt D}{2}U(D),
 \qquad D\ge10^{10}.
\]
This completes the short proof; no bound on $R_D$ was suppressed.

For completeness, the already inherited abelian-field formulation
NT-H04 is the following precise product calculation.  If $K/\mathbb Q$
is abelian with Dirichlet-character group $X_K$, then
\[
 h_KR_K=\frac{w_K\sqrt{\Delta_K}}{2^{r_1}(2\pi)^{r_2}}
              \prod_{\substack{\chi\in X_K\\\chi\ne1}}L(1,\chi).
\tag{NT-H04}
\]
Indeed Artin factorization gives the product for $\zeta_K/\zeta$,
and taking the residue at one gives this formula by the same class-number
formula.  Its left side and the prefactor are positive, so taking moduli
turns the product into $\prod|L(1,\chi)|$.  For any factors whose
conductors are at least $10^{10}$ insert $U$ and $1/V$; retain the finite
remaining product exactly or bound it separately.  This proves every
claimed product bound.  Dividing by an unestimated regulator does not
prove the corresponding bound on $h_K$.  This is the correction already
made in RETURN\_ZETA\_TWO IA02 and IA09, not a new result of this run.

\paragraph{NT-A01 and NT-A02: the nonabelian input in Artin's primitive-root problem.}
For a fixed integer $g\ne-1$ that is not a square, put
$K_n=\mathbb Q(\zeta_n,g^{1/n})$.  Hooley's classical theorem gives a
positive density
\[
 \delta(g)=\sum_{n\ge1}\frac{\mu(n)}{[K_n:\mathbb Q]}
\]
of primes where $g$ is a primitive root, assuming RH for these Kummer
fields.  We cite Hooley, \emph{On Artin's conjecture} (1967), the
unnumbered theorem on pp.~219--220.  The original internal page has not
been freshly obtained here; a fully inspected precise proof locator for
the density implication is Lenstra, \emph{On Artin's conjecture and
Euclid's algorithm in global fields} (1977), Theorem~(3.1), printed
p.~206, specialized to $K=F=\mathbb Q$, $W=\langle g\rangle$, $k=1$,
with density (2.15), p.~206, and positivity criterion (4.1), p.~207.
The series uses actual field degrees; replacing it by the universal
Artin constant for every $g$ loses the entanglement correction.

Lenstra's version also imposes a conjugacy-stable Frobenius condition
$C\subseteq\operatorname{Gal}(F/K)$ and asks that the residue image of
a finitely generated positive-rank $W\subseteq K^\times$ have index
dividing $k$.  Its auxiliary fields are formed from
$K(\zeta_{q(\ell)},W^{1/q(\ell)})$ with
$q(\ell)=\ell^{v_\ell(k)+1}$, taking the compositum for primes dividing
a squarefree $n$.  Theorem~(3.1) assumes RH for these fields' zeta
functions.  Theorem~(4.1) gives the finite obstruction criterion for
positive density.  $M$ supplies the hypotheses; $A$ alone does not
supply them because the Kummer fields need not be abelian over
$\mathbb Q$.  In geometric terms this is a tower of Kummer torsors of
$\mathbb G_m$, not merely the cyclotomic torsion cover.

\paragraph{NT-A03: a PID becomes Euclidean under the stated Kummer GRH.}
Lenstra's Theorem~(9.1), printed p.~218, states the result recorded in
NT-A03 for the ring of $S$-integers with $|S|\ge2$ and trivial ideal
class group.  Equation (9.2) gives its least Euclidean function.
The original page image confirms the weak inequality $|S|\ge2$;
the OCR's apparent strict inequality was not adopted.  The result is a
classical cited theorem with an exact proof locator, not a new proof of
the Euclidean algorithm.  The independent PID hypothesis and the
distinction from a norm-Euclidean algorithm remain in the catalogue.

\subsection{Growth, cancellation, and zero statistics}

\paragraph{NT-L02 and NT-L03. The fixed-character Lindel\"of deduction.}
The classical input is Simoni\v c, \emph{Estimates for L-functions in the
critical strip under GRH with effective applications} (2023), Theorem~1,
equations (4)--(6), proved in its \S3.  For a function $F$ in the Selberg
class with polynomial Euler product, zero-free to the right of the
critical line, it implies, for fixed $\delta>0$ and large $|t|$,
\[
 \log|F(\sigma+it)|,\qquad -\log|F(\sigma+it)|
 \ll_{F,\delta}(\log|t|)^{1-2\delta}\log\log|t|,
 \quad \tfrac12+\delta\le\sigma\le1.
\tag{NT-L-log}
\]
For $\delta\ge1/2$ one uses instead a smaller positive $\delta$ and
the absolutely convergent Euler product on $\sigma>1$.
The modulus bound for $\log F$ in the cited theorem gives both signs of
its real part.  Dividing the right side of (NT-L-log) by $\log|t|$
shows that it tends to zero for fixed $0<\delta<1/2$.  Hence, for every
$\epsilon>0$,
\[
 F(\sigma+it),\quad F(\sigma+it)^{-1}
 \ll_{F,\delta,\epsilon}(1+|t|)^\epsilon
 \qquad(\sigma\ge\tfrac12+\delta),
\tag{NT-L-right}
\]
with pole neighbourhoods omitted if $F$ has a pole.  The finite range of
$t$ is absorbed into the constant on any closed set avoiding poles and
zeros.  On $\sigma\ge1+\delta$ the Euler product gives the required
uniform bound directly; compact intermediate strips are already covered
by (NT-L-log) or the cited theorem's adjacent range.

Here is the passage to the critical line, often suppressed in invoking
Lindel\"of.  Fix $0<\delta<1/2$.  If $F$ has degree $d$, its functional
equation and Stirling's formula imply
\[
 |F(\tfrac12-\delta+it)|
 \ll_{F,\delta}(1+|t|)^{d\delta}
       |F^\vee(\tfrac12+\delta-it)|.
\]
The dual has the same zero-free hypothesis, so NT-L-right bounds this by
$O((1+|t|)^{d\delta+\eta})$ for every $\eta>0$.  The right boundary
$\sigma=1/2+\delta$ has exponent $\eta$.
The Phragm\'en--Lindel\"of principle on this strip gives at its midpoint
\[
 |F(\tfrac12+it)|\ll_{F,\delta,\eta}
                       (1+|t|)^{d\delta/2+\eta}.
\]
Its hypotheses hold because the standard functional equation and
finite-order analytic continuation give polynomial growth in each fixed
vertical strip; if there is a pole, apply the principle above and below
a fixed height containing it.  Choose $\delta$ and $\eta$ so that
$d\delta/2+\eta<\epsilon$.  The remaining bounded interval is compact,
and the critical line contains no pole for these L-functions.  This
proves $F(1/2+it)\ll_{F,\epsilon}(1+|t|)^\epsilon$.

For primitive Dirichlet functions the hypotheses are classical, and $A$
is the zero-free input; this proves NT-L02.  For a fixed non-pure-norm
unitary Hecke character, the Hecke L-function is entire of finite order,
its gamma shifts have nonnegative real parts, and its degree is
$[K:\mathbb Q]$.  On grouping its ideal Euler factors by rational primes,
one obtains a polynomial Euler product of that degree with local roots
of modulus at most one.  The Dirichlet coefficients are bounded by
$d_{[K:\mathbb Q]}(n)=O_\epsilon(n^\epsilon)$: each local coefficient
is bounded by the corresponding coefficient of $(1-p^{-s})^{-[K:\mathbb Q]}$.
Thus it meets the explicit axioms in Simoni\v c \S2.1, and $M$ supplies
zero-freeness.  If the character is a pure unitary norm power, its
L-function is a vertical translate of $\zeta_K$, to which the same
argument applies before translation.  Imprimitive Euler factors, when
included in the convention, form a fixed finite product; they and their
inverses are bounded on each $\sigma\ge1/2+\delta$, since their roots
have modulus at most one.  This proves NT-L03 in all the conventions just
specified.  Constants may depend on the field, character, infinity type,
and $\epsilon$; this proof makes no uniform conductor-aspect claim.

\paragraph{NT-L05 and NT-L06. Mertens, Liouville, and reciprocal tails.}
Titchmarsh, second edition, Theorem~14.25(C), printed p.~370, is the exact
classical equivalence
$\mathrm{RH}\Longleftrightarrow M(x)=O_\epsilon(x^{1/2+\epsilon})$
for every $\epsilon>0$.  We use that theorem rather than reproduce its
contour proof.  The following short deductions are included in full.
For a prime power $p^a$,
\[
 \sum_{d^2\mid p^a}\mu(p^a/d^2)=
 \begin{cases}1,&a\text{ even},\\-1,&a\text{ odd},\end{cases}
\]
because only the term with remaining exponent zero or one is nonzero.
Both sides are multiplicative, hence
$\lambda(n)=\sum_{d^2\mid n}\mu(n/d^2)$ for every $n$.
Consequently
\[
 \left|\sum_{n\le x}\lambda(n)\right|
 \le\sum_{d\le\sqrt x}|M(x/d^2)|
 \ll_\epsilon x^{1/2+\epsilon}
                \sum_{d\ge1}d^{-1-2\epsilon}
 \ll_\epsilon x^{1/2+\epsilon}.
\]
This proves Liouville cancellation, with the dependence on $\epsilon$
retained in the convergent sum.

Let $s=\sigma+it$ and choose $\eta>0$ with $1/2+\eta<\sigma$.
Partial summation for $y>x$ gives
\[
 \sum_{x<n\le y}\frac{\mu(n)}{n^s}
 =M(y)y^{-s}-M(x)x^{-s}
     +s\int_x^y M(u)u^{-s-1}\,du.
\]
The first term tends to zero as $y\to\infty$, and the integral converges
absolutely after using $|M(u)|\ll_\eta u^{1/2+\eta}$.  Its modulus is at
most a constant times
$|s|x^{1/2+\eta-\sigma}/(\sigma-1/2-\eta)$.
Together with the boundary term this proves the stated tail estimate.
It also proves local uniform convergence of the reciprocal Dirichlet
series on $\operatorname{Re}s>1/2$ by choosing $\eta$ below the distance
of a compact set from the boundary.  It agrees with $1/\zeta(s)$ on
$\operatorname{Re}s>1$ by the absolutely convergent Euler product and
therefore throughout that half-plane by analytic continuation.  At
$s=1$, $1/\zeta$ has value zero.  Hence $\sum\mu(n)/n=0$, as asserted.

\paragraph{NT-L04, NT-L07, NT-L08: what moments and zero density say.}
Harper (2013), Theorem~1, author p.~1, gives
$\int_T^{2T}|\zeta(1/2+it)|^{2k}dt\ll_kT(\log T)^{k^2}$ for fixed
$k\ge0$ under RH.  Its proof does not supply the conjectural leading
constant for every $k$.  Simoni\v c (2023), Theorem~2, equations (7),(8),
provides the explicit effective Mertens constants in NT-L07.
For NT-L08, if $\sigma>1/2$, every nontrivial zero covered by $A$ or
conditionally $M$ has real part $1/2<\sigma$, so the indicated counting
set is empty.  At $\sigma=1/2$ this reasoning says nothing about
multiplicity.  Trivial zeros from gamma factors and zeros of removed
finite Euler factors are not silently included in this definition.

\paragraph{NT-Z01 and NT-Z02: the exact pair-correlation range.}
For zeta zero ordinates $0<\gamma,\gamma'\le T$, counted with
multiplicity, put
\[
 F(\alpha,T)=\left(\frac{T}{2\pi}\log T\right)^{-1}
 \sum_{\gamma,\gamma'}T^{i\alpha(\gamma-\gamma')}
                    \frac4{4+(\gamma-\gamma')^2}.
\]
Montgomery (1973), \S1, unnumbered theorem, printed p.~181, proves under
RH the formula in NT-Z01 uniformly on each $0\le\alpha\le1-\epsilon$.
Corollary~2, p.~182, gives the proportion $2/3+o(1)$ of simple zeros.
The extension of the pair-correlation formula beyond that range is
explicitly a conjecture in that paper.  Adopting RH removes the theorem's
hypothesis; it does not replace its range by all $\alpha$.

\paragraph{NT-Z03--NT-Z05: low zeros of the prime-modulus Dirichlet family.}
For prime $q$ average over the $q-2$ nonprincipal characters, with scaled
ordinates $x_{\chi,j}=\gamma_{\chi,j}\log q/(2\pi)$.
Hughes--Rudnick, \emph{Linear statistics of low-lying zeros of
L-functions} (2003), Theorem~3.1, author p.~5, gives the unitary mean for
admissible even test functions with Fourier support in $[-2,2]$.
Theorem~4.1, p.~7, gives variance
$\int_{-1}^1|u|\widehat f(u)^2du$ with support in $[-1,1]$.
Theorem~5.1, p.~9, gives the Gaussian moments only for
$m<2/\alpha$ when the support is in $[-\alpha,\alpha]$.
These identities are unconditional with complex $\gamma$ allowed;
$A$ supplies the real-ordinate interpretation.  They are not an
unrestricted all-moment Gaussian theorem for one fixed test function.

In that source's ordering on p.~1, $\gamma_{\chi,1}$ is the first
\emph{nonnegative} ordinate, so a central zero is included with value zero.
With exactly this convention, Corollary~8.2, p.~20, gives
$\liminf_q\min_{\chi\ne\chi_0}x_{\chi,1}\le1/4$ under $A$.
Theorem~8.3, pp.~20--21, gives a positive proportion for
$x_{\chi,1}<\beta$ if $\beta\ge0.633$.
No strictly-positive-zero variant was inferred by discarding central
zeros.  The geometric comparison is with finite-field family monodromy
and its symmetry group; it does not identify the unrestricted
number-field limiting distribution.

\paragraph{NT-Z06 and NT-Z07: the additional gates.}
No argument in the cited sources proves full GUE, simplicity of every
zero, or central nonvanishing for every quadratic Dirichlet function
from zero location alone.  We record the specific missing correlation
or nonvanishing estimates, without claiming a logical independence
theorem.  For Iwaniec--Luo--Sarnak, \emph{Low lying zeros of families of
L-functions} (2000), the standing assumption on printed p.~56 explicitly
includes RH for the GL$_2$ functions in the family as well as Dirichlet
GRH.  Theorem~1.1, pp.~58--59, gives support $(-2,2)$ for fixed even
weight and squarefree level tending to infinity.  $A$ supplies its
Dirichlet part.  A universal Hecke statement does not identify a general
non-CM GL$_2$ newform with a rank-one Hecke character.  Thus the GL$_2$
hypothesis remains an additional input in this invocation.

\paragraph{NT-Z08 and NT-Z09: two exact classical equivalents.}
Rodgers--Tao, \texttt{arXiv:1801.05914v5}, Theorem~1, author p.~3,
proves $\Lambda\ge0$ for the classical de Bruijn--Newman constant.
The definition and RH equivalence in its \S1 give
$A\Rightarrow\Lambda\le0$, hence $\Lambda=0$.
For the deformation as normalized there, real-rootedness holds exactly
for $t\ge\Lambda$, yielding the two assertions in NT-Z08.
Lagarias, \texttt{arXiv:math/0008177v2}, Theorem~1.1, author p.~1,
is the exact locator for NT-Z09's divisor-sum inequality.
These are cited classical consequences, not additional new proofs of RH.

\subsection{Fixed-character cancellation and the local zero count}

\paragraph{NT-L09 and NT-L10. A full Perron transfer.}
Here is the elementary analytic form used for both rational and ideal
M\"obius sums.  Suppose
$G(s)=\sum_{n\ge1}a_nn^{-s}$ is absolutely convergent for
$\operatorname{Re}s>1$, $|a_n|\le d_D(n)$ for a fixed positive integer
$D$, and $G$ extends holomorphically to $\operatorname{Re}s>1/2$.
Assume that for every $\delta,\eta>0$,
\[
 G(\sigma+it)\ll_{\delta,\eta}(1+|t|)^\eta
 \quad\text{uniformly for }\tfrac12+\delta\le\sigma\le2.
\tag{NT-Perron-input}
\]
Then $\sum_{n\le x}a_n\ll_\epsilon x^{1/2+\epsilon}$ for every
$\epsilon>0$.

We prove the assertion, recording the truncation error as well as the
contour displacement.  It suffices to take $0<\epsilon<1/4$ and
half-integral $x\ge4$, because replacing any real $y$ by
$\lfloor y\rfloor+1/2$ leaves the partial sum unchanged and changes its
size by a bounded factor; the bounded remaining range is absorbed into
the constant.  Put
$c=1+1/\log x$, $T=x^2$, $\delta=\epsilon/4$ and
$\eta=\epsilon/8$.  The classical truncated Perron kernel gives
\[
 \sum_{n\le x}a_n=
 \frac1{2\pi i}\int_{c-iT}^{c+iT}G(s)\frac{x^s}{s}\,ds+R,
\]
with
\[
 |R|\ll\sum_{n\ge1}|a_n|(x/n)^c
           \min\left(1,\frac1{T|\log(x/n)|}\right).
\tag{NT-Perron-error}
\]
This is the usual Perron formula; an exact classical locator is
Titchmarsh, Lemma~3.12 and its proof, printed pp.~61--62.  The kernel
estimate follows by moving the integral of $y^s/s$ to the left or right
according to $y>1$ or $y<1$, taking the residue at zero in the first
case, and integrating the horizontal tails; the alternative trivial
bound handles $|\log y|\le T^{-1}$.

For $n\le x/2$ or $n\ge2x$ the denominator $|\log(x/n)|$ is bounded
below by $\log2$, so that part of (NT-Perron-error) is
\[
 \ll x^cT^{-1}\sum_{n\ge1}d_D(n)n^{-c}
 =x^cT^{-1}\zeta(c)^D\ll xT^{-1}(\log x)^D.
\]
For $x/2<n<2x$, use $d_D(n)\ll_{D,\eta}x^\eta$ and
$|\log(x/n)|\gg |x-n|/x$.  Since $x$ is half-integral,
$\sum_{x/2<n<2x}|x-n|^{-1}\ll\log x$.  Also $(x/n)^c\ll1$,
uniformly for these $x,n$.  This part of the error is therefore
$O(x^{1+\eta}T^{-1}\log x)$.  Both parts are smaller than
$x^{1/2+\epsilon}$.

Move the vertical line to $\sigma_0=1/2+\delta$.  The integrand has no
pole in the intervening rectangle: $G$ is holomorphic there and
$\sigma_0>0$.  On either horizontal segment,
(NT-Perron-input) bounds the integral by
$O(x^cT^{\eta-1})=O(x^{-1+2\eta})$.
On the new vertical segment it gives
\[
 \ll x^{\sigma_0}\left(1+\int_1^Tt^{\eta-1}\,dt\right)
 \ll_\eta x^{1/2+\delta}T^\eta
 =O(x^{1/2+\epsilon/2}).
\]
The bounded part $|t|\le1$ is harmless because $\sigma_0>0$ and $G$
is bounded on its compact segment.  This proves the asserted partial-sum
bound with room in the exponent.

Apply this lemma first to $G(s)=1/L(s,\chi)$.
Its coefficients are $a_n=\mu(n)\chi(n)$, including the usual zero
values at integers not coprime to the conductor.  Under $A$, the
reciprocal is holomorphic on $\operatorname{Re}s>1/2$; the principal
L-function's pole produces a zero of the reciprocal, not a pole.
NT-L02 supplies (NT-Perron-input), so the twisted bound follows.
Character orthogonality, for $(a,q)=1$, is the finite identity
\[
 \sum_{\substack{n\le x\\n\equiv a\ ({\rm mod}\ q)}}\mu(n)
 =\frac1{\varphi(q)}\sum_{\chi\ ({\rm mod}\ q)}
       \overline{\chi(a)}\sum_{n\le x}\mu(n)\chi(n).
\]
Imprimitive functions differ by finitely many Euler factors whose zeros
lie on $\operatorname{Re}s=0$, and their reciprocals are bounded in any
closed half-strip to the right of $1/2$.  Thus each term satisfies the
same fixed-character argument.  A finite sum proves the fixed-$q$
progression bound, with no uniformity in $q$ claimed.

For a unitary Hecke character, expand
\[
 \frac1{L(s,\eta)}
 =\sum_{(\mathfrak a,\mathfrak f)=1}
         \mu_K(\mathfrak a)\eta(\mathfrak a)(N\mathfrak a)^{-s}
 =\sum_{n\ge1}a_nn^{-s}.
\]
For each $n$, $|a_n|$ is at most the number of ideals of norm $n$.
If $[K:\mathbb Q]=D$, comparison of the rational-prime Euler factors
of $\zeta_K$ with $(1-p^{-s})^{-D}$ gives this number at most $d_D(n)$.
Indeed there are at most $D$ prime ideals above $p$, each of residue
degree at least one, and replacing every local factor by a degree-one
factor can only increase its nonnegative coefficients.  Holomorphy and
(NT-Perron-input) follow conditionally on $M$ from NT-L03, including
pure norm twists.  The lemma proves NT-L10.  Averaging over the finite
character group of any fixed ray class group then projects onto one ray
class exactly as above; all sums are over ideals prime to its modulus.
This completes every claimed cancellation statement.

\paragraph{NT-Z10. The fixed infinite-order Hecke zero count.}
Let $\eta$ be any fixed unitary Hecke character of a number field $K$.
Its number of nontrivial zeros, counted with multiplicity, in
\[
 0\le\operatorname{Re}s\le1,\qquad |\operatorname{Im}s-T|\le1
\]
is $O_{K,\eta}(\log(|T|+3))$, unconditionally.  This assertion includes
infinite-order characters; its constant is not asserted uniform in
their infinity types.

Here is a full local argument.  On any fixed vertical strip the finite
Hecke L-function has a polynomial upper bound in $|t|$, away from its
possible pole.  In the present argument this follows directly from the
unconditional first part of Simoni\v c's Theorem~1, equation (2), on
$\sigma\ge1/2$, and from the functional equation and Stirling's formula
on the reflected strip.  The Selberg-class hypotheses for a non-pure-norm
character were verified in NT-L03 without using GRH; a pure norm twist
is a translate of the Dedekind function.  Consequently constants $C,A>0$
exist such that, for large $|T|$,
\[
 \max_{|s-(2+iT)|\le4}|L(s,\eta)|\le C(|T|+3)^A.
\tag{NT-J-upper}
\]
For a pure norm twist the only pole has a fixed imaginary coordinate, so
choose $|T|$ sufficiently large that this disk does not contain it.
In all other cases the function is entire.

The Euler product at the centre gives a lower bound that a general
Schwartz Mellin transform need not have:
\[
 |L(2+iT,\eta)|
 \ge\prod_{\mathfrak p\nmid\mathfrak f}(1+(N\mathfrak p)^{-2})^{-1}
 \ge\prod_{\mathfrak p}(1-(N\mathfrak p)^{-2})
 =\zeta_K(2)^{-1}>0.
\tag{NT-J-lower}
\]
These comparisons follow term by term from $|\eta(\mathfrak p)|=1$
and $1/(1+u)\ge1-u$ for $0\le u<1$, and persist on taking limits of
the absolutely convergent products.

Let $n_T(3)$ count zeros in the disk of radius $3$ about $2+iT$.
Jensen's formula on the radius-$4$ disk, or its limiting-radius version
if a zero lies on the outer circle, implies
\[
 n_T(3)\log(4/3)
 \le\log\max_{|s-(2+iT)|\le4}|L(s,\eta)|
                -\log|L(2+iT,\eta)|.
\]
The right side is $O_{K,\eta}(\log(|T|+3))$ by (NT-J-upper) and
(NT-J-lower).  Every point with $0\le\operatorname{Re}s\le1$ and
$|\operatorname{Im}s-T|\le1$ is at distance at most $\sqrt5<3$
from the centre, so its zeros are included in $n_T(3)$.  The remaining
bounded interval of $T$ meets only finitely many zeros, by meromorphic
analyticity and the absence of an identically zero function; increasing
the constant covers it.  This proves the claim.

The finite-order conductor-explicit counterpart is
Greni\'e--Molteni2019, Appendix~A, Theorem~A.1, author p.~26.
NT-Z10 instead proves the fixed-character form needed for the local
arithmetic divisor, and supplies the infinite-order version of root
MT-M04.  The argument uses the lower bound at its Jensen centre and
polynomial vertical growth.  It does not repair a statement quantified
over every function in an arbitrary Schwartz Mellin-transform space.

\subsection{An explicit uniform cyclotomic specialization}

\paragraph{NT-P11 and NT-P12.}
For $q\ge3$ and $(a,q)=1$, define
\[
 \psi(x;q,a)=\sum_{\substack{p^m\le x\\p^m\equiv a\ ({\rm mod}\ q)}}\log p,
 \qquad
 \pi(x;q,a)=\#\{p\le x:p\equiv a\ ({\rm mod}\ q)\}.
\]
Put $e_q=1$ if $q\equiv2\pmod4$ and $e_q=0$ otherwise.
Then $A$ implies
\begin{align*}
 \left|\psi(x;q,a)-\frac{x}{\varphi(q)}\right|
 &\le\sqrt{x}\left[
 \left(\frac{\log x}{2\pi}+2\right)\log q+
 \frac{\log^2x}{8\pi}+2\right]+e_q\log x,\qquad x\ge1,\\
 \left|\pi(x;q,a)-\frac{\operatorname{Li}_2(x)}{\varphi(q)}\right|
 &\le\sqrt{x}\left[
 \left(\frac1{2\pi}+\frac3{\log x}\right)\log q+
 \frac{\log x}{8\pi}+\frac1{4\pi}+\frac6{\log x}\right]+e_q,
 \qquad x\ge2.
\end{align*}

For the full proof, take $L=\mathbb Q(\zeta_q)$ and $K=\mathbb Q$.
The Galois group is $(\mathbb Z/q\mathbb Z)^\times$, of order
$n_L=\varphi(q)$, and its conjugacy classes are singletons.
The cyclotomic discriminant formula gives
\[
 \Delta_L=q^{\varphi(q)}
    \prod_{p\mid q}p^{-\varphi(q)/(p-1)},\qquad
 \log\Delta_L\le\varphi(q)\log q.
\]
At every prime $p\nmid q$, arithmetic Frobenius acts on $\zeta_q$ by
$\zeta_q\mapsto\zeta_q^p$.  Thus its $m$th power is the element
represented by $p^m$ modulo $q$.  The use of arithmetic or geometric
Frobenius merely replaces the labelled class by its inverse; the estimate
is identical after that relabelling.  The Dirichlet factorization of
$\zeta_L$ makes RH for it a consequence of $A$.
Apply NT-C01 and NT-C02 with $\delta_C=1/\varphi(q)$ and insert the
discriminant upper bound; all displayed coefficients multiplying
$\log\Delta_L$ are positive on the stated ranges.

If $q$ is odd or $4\mid q$, the primes unramified in $L$ are exactly
those not dividing $q$, so the resulting sums are the progression sums
just defined.  If $q\equiv2\pmod4$, write $q=2m$ with $m$ odd.
Then $L=\mathbb Q(\zeta_m)$ and the prime $2$ is unramified, whereas it
is excluded from a reduced class modulo $q$.  For odd primes the class
condition modulo $m$ is equivalent to that modulo $q$, since reduction
$(\mathbb Z/q\mathbb Z)^\times\to(\mathbb Z/m\mathbb Z)^\times$ is
an isomorphism.  The unramified Chebotarev sum can therefore differ only
by the powers of $2$ whose Frobenius powers lie in the designated class.
Their total weight is at most
$\lfloor\log x/\log2\rfloor\log2\le\log x$.
For the prime-counting sum the discrepancy is at most one.  Applying the
triangle inequality proves both formulas, including this degeneration.

The small moduli are transparent separately:
\begin{align*}
 \psi(x;1,0)&=\psi(x),& \pi(x;1,0)&=\pi(x),\\
 \psi(x;2,1)&=\psi(x)-\lfloor\log x/\log2\rfloor\log2,&
 \pi(x;2,1)&=\pi(x)-\mathbf1_{x\ge2}.
\end{align*}
For $1\le x<2$ the floor is zero.  Hence NT-P01 and NT-P03 give the
corresponding bounds in their stated ranges, with the explicit powers-of-two
correction retained for modulus two.  These are estimates on the open
cyclotomic arithmetic curve, and the additional term is exactly its
removed-prime contribution.

\subsection{A specified quadratic family and central nonvanishing}

\paragraph{NT-Z11. The exact one-level-density input.}
Let $\mathcal D$ consist of the nonzero odd squarefree integers, of both
signs.  For $d\in\mathcal D$, $8d$ is a fundamental discriminant and
$\chi_{8d}$ is its primitive quadratic character.  Put
$\ell_X=\log(X/(2\pi e))$.  Let $w$ be an even, nonnegative, nonzero
Schwartz function and write
\[
 W_w(X)=\sum_{d\in\mathcal D}w(d/X),\qquad
 \mathcal Z_{w,\phi}(X)=\frac1{W_w(X)}
  \sum_{d\in\mathcal D}w(d/X)
       \sum_{\rho=1/2+i\gamma}\phi(\gamma\ell_X/(2\pi)),
\]
where the inner sum counts every nontrivial zero with multiplicity.
Use the Fourier convention
$\widehat\phi(u)=\int_{\mathbb R}\phi(t)e^{-2\pi iut}\,dt$.
For a real even Schwartz function $\phi$ with
$\operatorname{supp}\widehat\phi\subset(-2,2)$, $A$ implies
\[
 \lim_{X\to\infty}\mathcal Z_{w,\phi}(X)
 =\widehat\phi(0)-\frac12\int_{-1}^1\widehat\phi(u)\,du.
\tag{NT-Z11}
\]
This is the leading term of Fiorilli--Parks--S\"odergren,
\texttt{arXiv:1601.06833}, Theorem~1.1, author p.~2, with definitions
(1.1)--(1.2) on the same page.  Its proof uses Theorem~3.5, author p.~12, and is concluded on p.~14.
The GRH input in the prime-sum step is Dirichlet GRH
(Lemma~2.5, author p.~7), which $A$ supplies.  No unrestricted Fourier
support, individual central nonvanishing, or nonabelian GRH is being
asserted here.

\paragraph{NT-Z12. A full nonvanishing deduction, including deweighting.}
Under $A$, at least $15/16$ of this specified family have a nonzero
central value in the following precise sense:
\[
 \liminf_{X\to\infty}
 \frac{\#\{d\in\mathcal D:|d|\le X,\ L(1/2,\chi_{8d})\ne0\}}
      {\#\{d\in\mathcal D:|d|\le X\}}\ge\frac{15}{16}.
\tag{NT-Z12}
\]
The same lower bound holds for the weighted proportions defined by any
of the fixed weights $w$ in NT-Z11.  This is a deduction from the cited
classical density theorem, not a claim of a new nonvanishing theorem.
Here are all approximation and counting details.

Let $m_d$ be the order of vanishing at $1/2$.  The completed primitive
quadratic function has root number $+1$.  One way to see the sign without
choosing a Gauss-sum square root is to divide the symmetric completed
Dedekind functional equation of $\mathbb Q(\sqrt{8d})$ by that of
$\mathbb Q$: the factorization
$\zeta_{\mathbb Q(\sqrt{8d})}=\zeta L(s,\chi_{8d})$ and the real or
complex gamma factors identify the quotient, up to a nonzero constant,
with the completed quadratic function.  The gamma duplication formula
handles the imaginary quadratic case.  Thus this entire function is
even about $s=1/2$.  Its gamma factor is finite and nonzero at the
centre, so $m_d$ is an even nonnegative integer.  In particular
\[
 2\mathbf1_{L(1/2,\chi_{8d})=0}\le m_d.
\tag{NT-central-parity}
\]
These functional equations are the standard Hecke equations; a precise
source for the completed convention is Ramakrishnan--Valenza,
Theorem~7-19, printed p.~279.

Choose real even smooth functions $g_j$ compactly supported in $(-1,1)$
such that $\int g_j=1$ and
$g_j\to g=\tfrac12\mathbf1_{[-1,1]}$ in $L^2(\mathbb R)$.
For example, smooth even cutoffs of $[-1+1/j,1-1/j]$, normalized by
their integrals, have these properties.  Define
\[
 f_j(t)=\int g_j(u)e^{2\pi iut}\,du,\qquad
 \phi_j(t)=|f_j(t)|^2.
\]
Then $\phi_j$ is real, even, nonnegative and Schwartz,
$\phi_j(0)=1$, and
$\operatorname{supp}\widehat\phi_j\subset(-2,2)$.
The support assertion follows because
$\widehat\phi_j=g_j*\widetilde g_j$, where
$\widetilde g_j(u)=\overline{g_j(-u)}$.
Plancherel gives $f_j\to f$ in $L^2$, with
\[
 f(t)=\frac{\sin(2\pi t)}{2\pi t}.
\]
Moreover, Cauchy--Schwarz yields
\[
 \|\phi_j-|f|^2\|_1
 \le\|f_j-f\|_2(\|f_j\|_2+\|f\|_2)\longrightarrow0.
\]
Consequently their Fourier transforms converge uniformly.  Since
\[
 \widehat{|f|^2}(u)=
 \begin{cases}\tfrac12-\tfrac14|u|,&|u|\le2,\\0,&|u|>2,
 \end{cases}
\]
we obtain
\[
 \lim_{j\to\infty}
 \left(\widehat\phi_j(0)-\frac12
                   \int_{-1}^1\widehat\phi_j(u)\,du\right)
 =\frac12-\frac12\cdot\frac34=\frac18.
\tag{NT-test-limit}
\]
Under $A$ every ordinate in NT-Z11 is real.  The nonnegativity of
$\phi_j$ and $\phi_j(0)=1$ therefore imply
\[
 \frac1{W_w(X)}\sum_{d\in\mathcal D}w(d/X)m_d
 \le\mathcal Z_{w,\phi_j}(X).
\]
First take $\limsup_{X\to\infty}$ for a fixed $j$, use NT-Z11, and
then let $j\to\infty$.  The result and (NT-central-parity) give
\[
 \limsup_{X\to\infty}
 \frac{\sum_{d\in\mathcal D}w(d/X)
               \mathbf1_{L(1/2,\chi_{8d})=0}}{W_w(X)}\le\frac1{16}.
\tag{NT-weighted-vanishing}
\]
This proves the weighted assertion without applying the density theorem
directly to a nonsmooth Fourier transform or to support touching $2$.

For completeness, the ordinary count follows by a separate elementary
argument.  If $A_+(Y)$ counts positive odd squarefree integers up to $Y$,
the squarefree indicator gives
\begin{align*}
 A_+(Y)&=\sum_{\substack{k\le\sqrt Y\\k\ {\rm odd}}}\mu(k)
       \#\{n\le Y/k^2:n\ {\rm odd}\}\\
 &=\frac Y2\sum_{k\ {\rm odd}}\frac{\mu(k)}{k^2}+O(\sqrt Y)
 =\frac4{\pi^2}Y+O(\sqrt Y).
\end{align*}
Here the omitted tail is $O(Y\sum_{k>\sqrt Y}k^{-2})=O(\sqrt Y)$,
and the Euler product is
$\sum_{k\ {\rm odd}}\mu(k)/k^2=8/\pi^2$.
Thus $N(X):=\#\{d\in\mathcal D:|d|\le X\}
 =8X/\pi^2+O(\sqrt X)$.
Partial summation of this counting formula gives, for every fixed
smooth compactly supported $w$,
\[
 W_w(X)=\frac4{\pi^2}X\int_{\mathbb R}w(t)\,dt+O_w(\sqrt X).
\]
Indeed apply Stieltjes integration to the positive count $A_+$ and
its reflected negative count; the integrated error is bounded by a
constant depending on $\int_0^\infty\sqrt t\,|w'(t)|dt$.

Fix $\delta>0$.  Take such an even smooth $w$ equal to $1$ on
$[-1,1]$, between $0$ and $1$ everywhere, and supported in
$[-1-\delta,1+\delta]$.  The number of vanishing members with
$|d|\le X$ is bounded above by its weighted counterpart.  Dividing by
$N(X)$ and using (NT-weighted-vanishing) gives a $\limsup$ of at most
$\frac1{16}\cdot\frac12\int w\le(1+\delta)/16$.
Letting $\delta\downarrow0$ proves (NT-Z12).

The characters in this statement are quadratic rank-one Artin objects.
Their varying conductors describe a family of double covers of open
arithmetic curves.  The symplectic expression in NT-Z11 matches the
finite-field hyperelliptic-family comparison discussed immediately after
Theorem~1.1 of the cited paper.  This comparison does not enlarge the
analytic support range or turn the $15/16$ bound into nonvanishing for
each member.  The ordinary proportion proved here concerns the
$\chi_{8d}$ family exactly; it has not been silently relabelled as all
fundamental discriminants.

\subsection{Growing fields, class numbers, and restricted unramified towers}

\paragraph{NT-H05 and NT-H06. Exact growing-field inputs.}
For pairwise nonisomorphic number fields $K_i\ne\mathbb Q$, write
$g_i=\tfrac12\log\Delta_{K_i}$ and $n_i=[K_i:\mathbb Q]$.
Let $N_q(K_i)$ count prime ideals of norm $q$, where $q$ ranges over
rational prime powers.  The family is \emph{asymptotically exact} when
the following limits exist for every indicated index:
\[
 \phi_q=\lim_i\frac{N_q(K_i)}{g_i},\qquad
 \phi_{\mathbb R}=\lim_i\frac{r_1(K_i)}{g_i},\qquad
 \phi_{\mathbb C}=\lim_i\frac{r_2(K_i)}{g_i}.
\]
Let $h_i,R_i$ be the ordinary class number and full unit regulator, and
let $\kappa_i=\operatorname{Res}_{s=1}\zeta_{K_i}(s)$.
Tsfasman--Vl\u adu\c t, \texttt{arXiv:math/0205129},
GRH Theorem~7.2, author p.~44 (proof pp.~44--49), gives
\begin{align*}
 \lim_i\frac{\log(h_iR_i)}{g_i}
 &=1+\sum_q\phi_q\log\frac q{q-1}
      -\phi_{\mathbb R}\log2-\phi_{\mathbb C}\log(2\pi),\\
 \lim_i\frac{\log\kappa_i}{g_i}
 &=\sum_q\phi_q\log\frac q{q-1}.
\end{align*}
Only RH for the family of Dedekind functions is needed: $M$ supplies it;
$A$ suffices when every $K_i/\mathbb Q$ is abelian.  Neither normality
nor a tower is required in this theorem.  If $n_i/g_i\to0$, then
$N_q(K_i)\le n_i$ makes all the displayed invariants zero, so the
first limit is $1$.  This proves the usual Brauer--Siegel asymptotic
in this range without adding a normality hypothesis.

The same paper's GRH Theorem~3.1, author p.~19 (proof pp.~19--21),
supplies the basic inequality
\[
 \sum_q\frac{\phi_q\log q}{\sqrt q-1}
 +\alpha_{\mathbb R}\phi_{\mathbb R}
 +\alpha_{\mathbb C}\phi_{\mathbb C}\le1,
\tag{NT-basic}
\]
where
\[
 \alpha_{\mathbb R}=\tfrac12\log(8\pi)+\tfrac\pi4+\tfrac\gamma2,
 \qquad
 \alpha_{\mathbb C}=\log(8\pi)+\gamma,
\]
and $\gamma$ is Euler's constant.  Both original theorem pages were
inspected as images.  These are cited classical theorems.  The next
two consequences and their exact tower interpretations are proved below.

\paragraph{NT-H07. A root-discriminant bound with prescribed splitting.}
Suppose $n_i\to\infty$,
\[
 \Delta_{K_i}^{1/n_i}\longrightarrow\varrho\in(1,\infty),
 \qquad r_1(K_i)/n_i\longrightarrow\tau\in[0,1].
\]
Let $P$ be any set of rational primes such that every $p\in P$ splits
completely in every $K_i$.  Under RH for all $\zeta_{K_i}$,
\[
 \log\varrho\ge\log(8\pi)+\gamma+\frac\pi2\tau
          +2\sum_{p\in P}\frac{\log p}{\sqrt p-1}.
\tag{NT-H07}
\]
If $P$ is infinite, the inequality includes convergence of this
nonnegative series.  No normality assumption on $K_i/\mathbb Q$ is
needed.

To prove the claim, first take a subsequence with strictly increasing
degrees; its fields are pairwise nonisomorphic and its stated limits
are unchanged.  Note that
$g_i/n_i\to\tfrac12\log\varrho>0$.
For every fixed prime power $q$, the numbers $N_q(K_i)/g_i$ are
bounded since $N_q(K_i)\le n_i$.  Successive subsequence extraction,
followed by a diagonal subsequence over the countable prime powers,
makes all of them converge.  The archimedean ratios already converge;
on the resulting asymptotically exact subfamily they are
\[
 \phi_{\mathbb R}=\frac{2\tau}{\log\varrho},\qquad
 \phi_{\mathbb C}=\frac{1-\tau}{\log\varrho}.
\]
Complete splitting gives $N_p(K_i)=n_i$, hence
$\phi_p=2/\log\varrho$ for $p\in P$.
Apply (NT-basic) and discard the nonnegative terms from other norms.
For every finite $P_0\subset P$, multiplication by $\log\varrho$
yields
\[
 2\sum_{p\in P_0}\frac{\log p}{\sqrt p-1}
 +2\tau\alpha_{\mathbb R}+(1-\tau)\alpha_{\mathbb C}
 \le\log\varrho.
\]
The identity
$2\alpha_{\mathbb R}-\alpha_{\mathbb C}=\pi/2$
gives the claimed expression for finite $P_0$.  Taking the supremum
over finite subsets proves convergence and the same bound for all of
$P$.  The limiting quantities $\varrho,\tau$ are unchanged by the
subsequence, so this proves the original assertion.

\paragraph{NT-H08. Finiteness of a restricted arithmetic fundamental group.}
Fix a number field $K$ and a finite set $P$ of rational primes that
split completely in $K$.  Let $K^{\rm ur,P}$ be the compositum of all
finite extensions of $K$ unramified at every finite prime, in which
every prime of $K$ above $P$ splits completely.  Put
\[
 C(P)=8\pi\exp\left(\gamma+
                 2\sum_{p\in P}\frac{\log p}{\sqrt p-1}\right).
\]
If $\Delta_K^{1/[K:\mathbb Q]}<C(P)$, then $M$ implies
$[K^{\rm ur,P}:K]<\infty$.
If $K$ is totally real and only totally real extensions are allowed,
the corresponding conclusion holds under the weaker inequality
$\Delta_K^{1/[K:\mathbb Q]}<e^{\pi/2}C(P)$.
The inequalities are strict: no equality case is claimed.

Here is the full proof.  Being unramified at finite primes and splitting
completely at specified primes are both preserved by conjugation over
$K$ and by finite composita.  Thus the displayed compositum is Galois
over $K$, and all of its finite subextensions have these properties.
If it had infinite degree, it would contain a sequence of finite
Galois subextensions $K_i/K$ with degrees tending to infinity.
The relative discriminant formula and finite unramifiedness give
\[
 \Delta_{K_i}=\Delta_K^{[K_i:K]},\qquad
 \Delta_{K_i}^{1/[K_i:\mathbb Q]}=
                  \Delta_K^{1/[K:\mathbb Q]}=: \varrho.
\]
If $K=\mathbb Q$, this equality gives discriminant one for any such
$K_i$; Minkowski's discriminant bound already rules out a nontrivial
extension.  Otherwise $\varrho>1$.
After taking a subsequence, $r_1(K_i)/[K_i:\mathbb Q]$ tends to some
$\tau\in[0,1]$.  Every $p\in P$ splits completely in $K_i/\mathbb Q$
because it splits in $K$ and each prime above it splits in $K_i/K$.
$M$ supplies the Dedekind RH hypotheses of NT-H07.  That result yields
$\varrho\ge e^{\pi\tau/2}C(P)\ge C(P)$, a contradiction.
In the totally real version all $K_i$ are totally real, so $\tau=1$
and the same proof gives the second threshold.  This completes the proof.

For $P=\varnothing$, the thresholds are
$8\pi e^\gamma\simeq44.763232$ and
$8\pi e^{\gamma+\pi/2}\simeq215.332516$.
The group $\operatorname{Gal}(K^{\rm ur,P}/K)$ is the quotient of
$\pi_1^{\text{ét}}(\operatorname{Spec}\mathcal O_K)$ obtained by
killing the closed normal subgroup generated by Frobenius elements at
the primes above $P$.  Requiring total reality also kills the
archimedean complex-conjugation classes.  This is the finite-\'{e}tale
Galois-category dictionary, SGA~1, Expos\'e~V, \S7.  The discriminant
criterion therefore gives a concrete conditional finiteness statement
for this quotient.  It is a restriction on arithmetic towers in the
cote-149/GT comparison programme, not a claim of a new anabelian
reconstruction theorem.

\subsection{Prime races: zero location versus relations among ordinates}

\paragraph{NT-R01 and NT-R02. What $A$ alone supplies.}
Fix distinct reduced classes $a_1,\ldots,a_r$ modulo $q$, and set
\[
 E_{q;\boldsymbol a}(x)=\frac{\log x}{\sqrt x}
 \bigl(\varphi(q)\pi(x;q,a_1)-\pi(x),\ldots,
       \varphi(q)\pi(x;q,a_r)-\pi(x)\bigr).
\]
Rubinstein--Sarnak, \emph{Chebyshev's Bias} (1994), Theorem~1.1,
printed p.~174, proof \S2.1, gives under $A$ a probability measure
$\mu_{q;\boldsymbol a}$ such that, for every bounded continuous $f$,
\[
 \lim_{X\to\infty}\frac1{\log X}
       \int_2^X f(E_{q;\boldsymbol a}(x))\frac{dx}{x}
   =\int_{\mathbb R^r}f\,d\mu_{q;\boldsymbol a}.
\tag{NT-R01}
\]
Its normalization uses $\pi(x)$ in each coordinate.  Theorem~1.2,
printed p.~175, proof \S2.2, and the discussion following it imply
\[
 \underline\delta\{x:\pi(x;4,3)>\pi(x;4,1)\}>0,
 \qquad
 \underline\delta\{x:\pi(x;4,1)>\pi(x;4,3)\}>0,
\tag{NT-R02}
\]
where $\underline\delta S=\liminf_{X\to\infty}
(\log X)^{-1}\int_{S\cap[2,X]}dx/x$.
The same conclusion holds for the two classes modulo three.
These statements concern logarithmic sampling.  The existence of the
limiting distribution for continuous tests does not by itself justify
substituting an ordering-cone indicator: its boundary could have mass.

\paragraph{NT-R03 and NT-R04. The separate ordinate-independence input.}
Use the source's Grand Simplicity Hypothesis (GSH): the nonnegative
ordinates of zeros of all primitive Dirichlet functions, listed with
multiplicity, are linearly independent over $\mathbb Q$.  In particular
this excludes central zeros and repeated ordinates.  GSH remains an
additional hypothesis here.

Under $A$ and GSH, for $r\ge2$ every strict ordering race has a
logarithmic density in $(0,1)$; the source's formula (1.2) and
Remark~1.3, printed p.~176, give the requisite regularity and positive
mass.  For fixed $r$,
\[
 \max_{\substack{a_1,\ldots,a_r\pmod q\\
                   \text{distinct reduced classes}}}
 \left|\delta\{x:\pi(x;q,a_1)>\cdots>\pi(x;q,a_r)\}
                   -\frac1{r!}\right|\longrightarrow0
       \quad(q\to\infty).
\tag{NT-R04}
\]
This is Theorem~1.5, printed p.~177, proof \S3.2.
For the exact symmetry classification, put
$c(q,a)=-1+\#\{b\pmod q:b^2\equiv a\pmod q\}$.
Theorem~1.4, printed p.~177, proof Proposition~3.1, says that the
limiting law is invariant under every coordinate permutation precisely
in the following cases: $r=2$ and $c(q,a_1)=c(q,a_2)$; or $r=3$
and, after labelling, the classes are
$a_1,a_1\rho,a_1\rho^2$ with $\rho\ne1$ and $\rho^3=1\pmod q$.
This is symmetry of the full law.  It must not be conflated with a
converse assertion about one numerical race density.  The discussion
on p.~177 gives, in particular, the strict bias
$1/2<\delta\{x:\pi(x;4,3)>\pi(x;4,1)\}<1$ under these extra
hypotheses.  The relevant original pp.~174--178 were visually inspected
because their PDF text encoding is defective.

Geometrically these are races between Frobenius classes in cyclotomic
covers of $\operatorname{Spec}\mathbb Z[1/q]$.  Their limiting
distribution is an explicit-formula statistic of the corresponding
rank-one objects.  The additional GSH gate concerns relations among
zero ordinates, so the zero-location inputs $A$ and $M$ do not discharge
it in this argument.  A finite-field monodromy comparison likewise
does not remove that number-field hypothesis.

\subsection{Multiplicity, gaps, and the first zero: quantitative limits}

\paragraph{NT-Z13. An argument bound with its actual archimedean hypothesis.}
Carneiro--Chandee--Milinovich, \emph{A note on the zeros of zeta and
$L$-functions} (2015), author PDF Theorem~1, p.~2, proves under RH
\[
 |S(t)|\le \frac14\frac{\log t}{\log\log t}
       +O\left(\frac{\log t\log\log\log t}{(\log\log t)^2}\right).
\tag{NT-Z13a}
\]
Here $S(t)=\pi^{-1}\arg\zeta(1/2+it)$, with the argument continued
from $2$ through $2+it$ and the mean of its two one-sided values used
at a zero.  The proof is in \S3, pp.~3--5.

We record the precise Hecke specialization of their Theorem~5,
p.~7, equations (4.4)--(4.5), whose proof occupies pp.~7--12.
Let $\eta$ be a primitive finite-order Hecke character of a number
field $K$ of degree $n$.  Write its completion, up to a nonzero
constant independent of $s$, as
\[
 \Lambda(s,\eta)=G_\eta(s)L(s,\eta),\qquad
 G_\eta(s)=Q_\eta^{s/2}\prod_{j=1}^n\Gamma_{\mathbb R}(s+\mu_j),
 \qquad Q_\eta=\Delta_K\operatorname N\mathfrak f_\eta,
\]
where $\Gamma_{\mathbb R}(s)=\pi^{-s/2}\Gamma(s/2)$.
Each real place contributes $\mu_j=0$ or $1$ according to the sign
character.  Each complex place contributes the pair $0,1$.
Set
\[
 C_\eta=Q_\eta\prod_{j=1}^n(\mu_j+3),\quad
 C_\eta(t)=C_\eta(|t|+1)^n,\quad
 D_\eta(t)=\log\log\bigl(C_\eta(t)^{3/n}\bigr).
\]
For $t>0$ define
$S_\eta(t)=-\pi^{-1}\int_{1/2}^{\infty}
\operatorname{Im}(L'/L)(\sigma+it,\eta)\,d\sigma$ away from zero
ordinates, and use one-sided means at them.  Under RH for this
particular $L$-function,
\[
 |S_\eta(t)|\le\frac{\log C_\eta(t)}{4D_\eta(t)}
   +O\left(\frac{\log C_\eta(t)\log D_\eta(t)}{D_\eta(t)^2}\right).
\tag{NT-Z13b}
\]
The implied constant is absolute.  In particular, for fixed $K,\eta$,
\[
 |S_\eta(t)|\le\frac n4\frac{\log t}{\log\log t}
   +O_{K,\eta}\left(
       \frac{\log t\log\log\log t}{(\log\log t)^2}\right).
\tag{NT-Z13c}
\]

Here is the check of the source hypotheses, which matters for $M$.
After grouping prime ideals by their underlying rational prime, the
Hecke Euler factor is a product of at most $n$ linear reciprocal
factors with local parameters of modulus at most one: factor each
$1-\eta(\mathfrak p)p^{-fs}$ into its $f$ roots and pad missing
factors by zero parameters.  Thus the source's parameter $\vartheta$
is zero.  The completion has order one, the standard functional
equation with the conjugate character, and is entire except for the
trivial character's poles at $0,1$.  Its listed gamma parameters are
real and nonnegative.  These are precisely the analytic conditions
in \S4, p.~6, supplied by classical Hecke theory (RV1999,
Theorem~7-19).  $M$ supplies the remaining zero-location assumption;
$A$ suffices for Dirichlet characters over $\mathbb Q$.
In particular, the source explicitly requires gamma parameters to be
real or to occur in conjugate pairs.  We do not silently use it for
an arbitrary infinite-order unitary Hecke character.  The paired
$\pm it_k$ family in root MT-H01--H03 can be checked separately
against this same hypothesis.

\paragraph{NT-Z14. Full deduction of multiplicity and gap bounds.}
For each fixed $K,\eta$ as in NT-Z13, under its GRH, a zero
$1/2+i\gamma$ with $\gamma\to+\infty$ has multiplicity
\[
 m_\eta(\gamma)\le\frac n2\frac{\log\gamma}{\log\log\gamma}
   +O_{K,\eta}\left(
    \frac{\log\gamma\log\log\log\gamma}{(\log\log\gamma)^2}\right).
\tag{NT-mult}
\]
For consecutive distinct sufficiently large positive ordinates
$\gamma<\gamma'$,
\[
 \gamma'-\gamma\le\frac{\pi}{\log\log\gamma}
   +O_{K,\eta}\left(
    \frac{\log\log\log\gamma}{(\log\log\gamma)^2}\right).
\tag{NT-gap}
\]
Consequently every sufficiently large real $T$ lies within
\[
 \frac{\pi}{2\log\log T}
 +O_{K,\eta}\left(\frac{\log\log\log T}{(\log\log T)^2}\right)
\]
of a zero ordinate.  The analogous
negative-height statements follow by conjugating $\eta$.  These are
classical consequences in a fixed-character formulation; they assert
neither simplicity nor pair correlation.

We give the entire deduction, including the jump and endpoint issues.
In a sufficiently high horizontal region $L(s,\eta)$ has no pole;
GRH and the boundary nonvanishing theorem give no zero to the right
of the critical line.  Continue $\log L$ from its Euler-product
normalization there.  Near $\rho=1/2+i\gamma$ write
$L(s,\eta)=(s-\rho)^m U(s)$ with $U$ holomorphic and nonzero.
Approaching the critical line from its right, the argument of
$s-\rho$ changes from $-\pi/2$ to $\pi/2$ across the ordinate.
Hence
\[
 S_\eta(\gamma+)-S_\eta(\gamma-)=m.
\]
Taking one-sided limits in NT-Z13c bounds this jump by twice its
right-hand side and proves (NT-mult).

Let $\Theta_\eta(t)=\operatorname{Im}\log G_\eta(1/2+it)$, with
any continuous determination at large positive $t$.
The functional equation and conjugation give
$\Lambda(1/2+it,\eta)=\varepsilon_\eta
\overline{\Lambda(1/2+it,\eta)}$.
Therefore, on an interval containing no zero, the argument of the
completed function is constant modulo $\pi$, and
\[
 S_\eta'(t)=-\frac1\pi\Theta_\eta'(t),\qquad
 \frac1\pi\Theta_\eta'(t)
    =\frac n{2\pi}\log t+O_{K,\eta}(1).
\tag{NT-phase}
\]
The derivative formula is Stirling's formula applied to each of the
$n$ displayed $\Gamma_{\mathbb R}$ factors; the conductor term is
constant.  The same identity and NT-Z13c show that positive ordinates
are unbounded: if there were no zeros beyond a height, integration
would force $S_\eta(t)$ to have size comparable to $t\log t$.

For large $\gamma$, a zero-free gap cannot have length at least one.
Otherwise it contains an interval of length $1/2$ whose endpoints
are between $\gamma$ and $\gamma+1$.  Integrating (NT-phase) over it
gives a decrease of size $(n/(4\pi)+o(1))\log\gamma$, whereas
NT-Z13c bounds the difference of its endpoint values by
$O_{K,\eta}(\log\gamma/\log\log\gamma)$, a contradiction.
Thus $h=\gamma'-\gamma<1$ eventually.  Integrate (NT-phase) from
just above $\gamma$ to just below $\gamma'$ and then take endpoint
limits.  Uniformly on this interval,
\[
 \left(\frac n{2\pi}\log\gamma+O_{K,\eta}(1)\right)h
   =S_\eta(\gamma+)-S_\eta(\gamma'-)
   \le\frac n2\frac{\log\gamma}{\log\log\gamma}
      +O_{K,\eta}\left(
       \frac{\log\gamma\log\log\log\gamma}
            {(\log\log\gamma)^2}\right).
\]
Division proves (NT-gap).  Any $T$ between consecutive ordinates is
at distance at most half their gap from one endpoint.  Since that
gap is less than one, replacing $\gamma$ by $T$ in the slowly varying
bound changes only the error term.  This proves the last assertion
and completes the proof.  For $K=\mathbb Q$ and the trivial
character, $n=1$ gives the usual zeta multiplicity and gap bounds
under $A$ alone.

\paragraph{NT-Z15 and NT-Z16. Conductor aspect and the first ordinate.}
In the same finite-order Hecke setting put
$D_\eta=\log\log(C_\eta^{3/n})$.
Theorem~5, equation (4.5), gives
\[
 \operatorname{ord}_{s=1/2}L(s,\eta)
 \le\frac{\log C_\eta}{2D_\eta}
   +O\left(\frac{\log C_\eta\log D_\eta}{D_\eta^2}\right).
\tag{NT-Z15}
\]
There is no gamma pole or zero at $1/2$, so the order equals that
of the completion in the source.  Theorem~7, p.~13, with its full
proof on that page, gives at least one zero $1/2+i\gamma$ satisfying
\[
 |\gamma|\le\frac{\pi}{2D_\eta}
          +O\left(\frac{\log D_\eta}{D_\eta^2}\right)
 \quad\text{as } C_\eta^{3/n}\longrightarrow\infty.
\tag{NT-Z16}
\]
Both implied constants are universal.  A possible central zero is
included: this is not a bound for the first strictly positive
ordinate.  Keeping $n$ inside $D_\eta$ is essential when the degree
varies.  For primitive Dirichlet characters of conductor $q>2$, the
specializations are
\[
 \operatorname{ord}_{1/2}L(s,\chi)
       \le(1/2+o(1))\frac{\log q}{\log\log q},\qquad
 \min_\rho|\operatorname{Im}\rho|
       \le\frac{\pi/2+o(1)}{\log\log q}.
\]
Here $A$ supplies every required GRH hypothesis.  For finite-order
Hecke characters over general fields the corresponding input is
the finite-order part of $M$; it is not supplied by $A$ merely
because the base field happens to be abelian over $\mathbb Q$.

These bounds constrain the zero divisor of rank-one arithmetic
local systems and their Artin inductions.  In the Grothendieck
comparison language they are quantitative restrictions on the
archimedean explicit-formula spectrum.  They supply neither a
finite-field purity argument over $\operatorname{Spec}\mathbb Z$
nor a general geometric multiplicity-one theorem.  They refine the
specific open gates in NT-Z06 while leaving them open.

\subsection{Further prime constants and quadratic-family moments}

\paragraph{NT-P13. Explicit and asymptotic prime-gap constants.}
Under the RH part of $A$, every $x\ge4$ has a prime in
\[
 \left[x,x+\frac{22}{25}\sqrt x\log x\right].
\]
This is Carneiro--Milinovich--Soundararajan,
\emph{Fourier optimization and prime gaps} (2019), author v2
Theorem~5, p.~5; the explicit proof is \S6, pp.~21--26.
Consequently, for every consecutive pair with $p_n>3$,
$p_{n+1}-p_n\le(22/25)\sqrt{p_n}\log p_n$.
For clarity about the closed left endpoint, apply the interval
statement to $x=p_n+\varepsilon$ with
$0<\varepsilon<p_{n+1}-p_n$, then let $\varepsilon\downarrow0$.
Corollary~4 on the same page, following Theorem~3 and its proof in
\S5, strengthens the asymptotic constant:
\[
 \limsup_{n\to\infty}
       \frac{p_{n+1}-p_n}{\sqrt{p_n}\log p_n}<\frac{21}{25}.
\]
The explicit all-$x$ statement and the sharper asymptotic statement
have different quantifiers.  No claim is made that these are the
latest possible constants.

\paragraph{NT-P14--NT-P16. Fourier improvements in the conductor aspect.}
Carneiro--Milinovich--Quesada-Herrera--Ramos,
\emph{Fourier optimization, the least quadratic non-residue, and
the least prime in an arithmetic progression}, arXiv v2 dated
August 2025, Theorems~1 and 2, author p.~2, gives under $A$
\[
 \limsup_{p\to\infty}\frac{n_p}{(\log p)^2}<0.7615,
 \qquad
 \limsup_{p\to\infty}\frac{r_p}{(\log p)^2}<0.7615.
\tag{NT-P14}
\]
Here $p$ runs over primes, $n_p$ is the least positive quadratic
nonresidue modulo $p$, and $r_p$ is the least prime $r$ with
$(r/p)=1$.  The nonresidue is itself prime.  Theorem~1 follows from
Theorem~6, whose proof is \S4, pp.~13--17; the modification proving
Theorem~2 is \S4.4, pp.~16--17.  These bounds improve the
asymptotic constant; they do not replace NT-P07's explicit threshold
$p\ge5$ by an unstated numerical threshold.

More generally, for a nonprincipal character $\chi$ modulo $q$ of
fixed order $\ell\ge2$, define
$n_\chi=\min\{n\ge1:\chi(n)\notin\{0,1\}\}$.
Theorem~6, p.~5, and Corollary~7, p.~6, give
\[
 n_\chi\le(c_\ell+o_\ell(1))(\log q)^2,
 \qquad
 \begin{array}{c|rrrr}
   \ell&2&3&4&5\\ \hline
   c_\ell&0.7615&0.6707&0.6131&0.5765
 \end{array}
\tag{NT-P15}
\]
as $q\to\infty$.  For fixed $\ell\ge6$ the same corollary permits
\[
 c_\ell=\frac14\left(
  1-\frac{\frac{\ell}{\ell-1}\log\bigl((3\ell-4)/\ell\bigr)}
             {\log(\ell-1)}+\frac1{\ell-1}\right)^{-2}.
\]
The character need not be primitive.  Its values zero at primes
dividing $q$ are explicitly excluded in the definition.  If
$n_\chi$ were composite, every prime divisor would be smaller and
coprime to $q$, hence would have character value one by minimality;
multiplicativity would give $\chi(n_\chi)=1$, a contradiction.
Thus these are also bounds for the least prime outside the kernel
of the cyclotomic character.  The limit $c_\ell\to1/4$ is a
statement about these constants; we do not assert uniformity when
$\ell$ varies with $q$.

For the least prime $P(a,q)\equiv a\pmod q$, the same paper's
Theorem~3, author p.~3, proof \S5, pp.~18--22, gives under $A$
\[
 \limsup_{q\to\infty}
  \max_{(a,q)=1}\frac{P(a,q)}{(\varphi(q)\log q)^2}<\frac89.
\tag{NT-P16}
\]
The uniformity in the reduced class follows from the proof: its
errors in (5.15)--(5.16) depend on the fixed Fourier test function
and not on $a$.  The numerical optimization certification in the
paper is cited, not rerun here.  These results concern closed points
and selected Frobenius classes of cyclotomic covers.  They sharpen
the effective detection bounds in the SGA~1 Galois-category
dictionary; they do not enlarge the class of $L$-functions furnished
by $A$.

\paragraph{NT-L11. Quadratic central moments with the full family specified.}
For each fixed real $k\ge0$, $A$ and Harper,
\emph{Sharp conditional bounds for moments of the Riemann zeta
function}, Theorem~2, author p.~15, give
\[
 \sum_{\substack{X\le |d|\le2X\\d\ {
m fundamental}}}
       |L(1/2,\chi_d)|^k
     \ll_k X(\log X)^{k(k+1)/2}.
\tag{NT-L11}
\]
The family includes positive and negative fundamental discriminants;
$\chi_d$ is the associated primitive quadratic character.  The
argument is in \S5, pp.~15--17, using Propositions~3 and 4.
The input is only GRH for the quadratic functions in question.
This is an upper bound of the predicted order, not an asymptotic
formula with its leading constant.  In particular it asserts no
bound for negative moments or reciprocals of central values.

\paragraph{NT-L12. Full fixed-threshold large-value deductions.}
For every fixed $\alpha>0$, $A$ implies
\[
 \operatorname{meas}\{t\in[T,2T]:
            |\zeta(1/2+it)|\ge(\log T)^\alpha\}
       \ll_\alpha \frac{T}{(\log T)^{\alpha^2}}.
\tag{NT-L12a}
\]
For every fixed $\alpha>1/2$ it also implies
\[
 \#\{d\text{ fundamental}:X\le|d|\le2X,
          |L(1/2,\chi_d)|\ge(\log X)^\alpha\}
       \ll_\alpha
          \frac{X}{(\log X)^{(\alpha-1/2)^2/2}}.
\tag{NT-L12b}
\]
Here is the complete proof.  For any fixed $k>0$, Markov's
inequality and Harper's Theorem~1 (NT-L04) bound the first measure by
\[
 (\log T)^{-2k\alpha}
       \int_T^{2T}|\zeta(1/2+it)|^{2k}\,dt
       \ll_k T(\log T)^{k^2-2k\alpha}.
\]
Choose the fixed value $k=\alpha$ to obtain (NT-L12a).
Similarly NT-L11 bounds the second count by
$O_k(X(\log X)^{k(k+1)/2-k\alpha})$.
For $\alpha>1/2$, choose the positive fixed value
$k=\alpha-1/2$; completing the square gives
$k(k+1)/2-k\alpha=-(\alpha-1/2)^2/2$.
This proves (NT-L12b).  All constants may depend on the fixed
threshold exponent; no choice of $k$ growing with $T$ or $X$ has
been made.  Central zeros cause no difficulty because only positive
moments occur.  The result is a direct optimization of the cited
moment bounds, not a new moment theorem.

The quadratic family consists of rank-one Artin objects coming from
quadratic covers with growing discriminant.  Its central-value
moment exponent is an arithmetic family statement.  A geometric
symmetry-type comparison is useful context, but is not substituted
for Harper's analytic proof or the exact GRH input.

\subsection{Closed-point Frobenius relations in unramified towers}

\paragraph{NT-H09. A relative splitting-degree budget.}
Let $K\ne\mathbb Q$ be a number field of degree $n$, discriminant
$\Delta_K$, and root discriminant $\varrho=\Delta_K^{1/n}$.
Let $S$ be a finite set of prime ideals of $K$, with
$q_{\mathfrak p}=\operatorname N\mathfrak p$, and assign to each
$\mathfrak p\in S$ a positive integer $m_{\mathfrak p}$.
Suppose that $K_i/K$ are finite Galois extensions unramified at all
finite primes, $d_i=[K_i:K]\to\infty$, and the common residue
degree of primes above $\mathfrak p$ in $K_i/K$ is at most
$m_{\mathfrak p}$ for every $i$.  Assume
$r_1(K_i)/[K_i:\mathbb Q]\to\tau$ and RH for every
$\zeta_{K_i}$.  Then
\[
 \log\varrho\ge\log(8\pi)+\gamma+\frac\pi2\tau
  +\frac2n\sum_{\mathfrak p\in S}
       \frac{\log q_{\mathfrak p}}
            {q_{\mathfrak p}^{m_{\mathfrak p}/2}-1}.
\tag{NT-H09}
\]
Thus $M$ supplies the needed zero hypotheses.  No nesting of the
extensions is required.  Complete splitting is the case
$m_{\mathfrak p}=1$ and does not require the rational prime below
$\mathfrak p$ to split in $K$.

Here is the full proof.  Pass to strictly increasing degrees so that
the fields are pairwise nonisomorphic.  Finite unramifiedness gives
$\Delta_{K_i}=\Delta_K^{d_i}$ and
$g_i=\tfrac12 d_i\log\Delta_K$.
There are finitely many possible residue degrees at each point of
the finite set $S$, so pass to a subsequence on which these degrees
are fixed, say $f_{\mathfrak p}\le m_{\mathfrak p}$.
As in NT-H07, boundedness of $N_Q(K_i)/g_i$ and a diagonal
subsequence over rational prime powers $Q$ give an asymptotically
exact family.  Neither the limiting signature nor the displayed
residue degrees change.  Since the extension is Galois and
unramified, precisely $d_i/f_{\mathfrak p}$ primes lie above
$\mathfrak p$, each of norm $q_{\mathfrak p}^{f_{\mathfrak p}}$.
Different base prime ideals give disjoint sets of primes upstairs,
even when their absolute norms coincide.  Therefore
\[
 \phi_Q\ge\frac2{\log\Delta_K}
       \sum_{\substack{\mathfrak p\in S\\
                    q_{\mathfrak p}^{f_{\mathfrak p}}=Q}}
           \frac1{f_{\mathfrak p}}.
\]
The archimedean invariants are
$\phi_{\mathbb R}=2n\tau/\log\Delta_K$ and
$\phi_{\mathbb C}=n(1-\tau)/\log\Delta_K$.
Substitute these lower bounds in the TV inequality (NT-basic),
discard the other nonnegative terms, and multiply by
$\log\Delta_K/n$.  Since
$\log(q_{\mathfrak p}^{f_{\mathfrak p}})/f_{\mathfrak p}
=\log q_{\mathfrak p}$, the finite-place contribution is at least
\[
 \frac2n\sum_{\mathfrak p\in S}
      \frac{\log q_{\mathfrak p}}
           {q_{\mathfrak p}^{f_{\mathfrak p}/2}-1}
 \ge\frac2n\sum_{\mathfrak p\in S}
      \frac{\log q_{\mathfrak p}}
           {q_{\mathfrak p}^{m_{\mathfrak p}/2}-1}.
\]
The archimedean part simplifies using
$2\alpha_{\mathbb R}-\alpha_{\mathbb C}=\pi/2$ exactly as in
NT-H07.  This proves (NT-H09), including the possibility of
coincident prime-power norms.

\paragraph{NT-H10. A finite quotient from bounded-order Frobenius relations.}
Let
$G_K=\pi_1^{\text{ét}}(\operatorname{Spec}\mathcal O_K)$, and
for the data $S,\boldsymbol m$ above let $N$ be the closed normal
subgroup generated by
$\operatorname{Frob}_{\mathfrak p}^{m_{\mathfrak p}}$ for
$\mathfrak p\in S$.  The normal subgroup is independent of the
choices inside the Frobenius conjugacy classes.  Put
\[
 C_K(S,\boldsymbol m)=8\pi\exp\left(
       \gamma+\frac2n\sum_{\mathfrak p\in S}
       \frac{\log q_{\mathfrak p}}
            {q_{\mathfrak p}^{m_{\mathfrak p}/2}-1}\right).
\]
Under $M$, if $\varrho<C_K(S,\boldsymbol m)$, then $G_K/N$ is
finite; equivalently, $N$ is open.  For totally real $K$, first
quotient $G_K$ by the closed normal subgroup generated by complex
conjugations.  In that totally real quotient the weaker condition
$\varrho<e^{\pi/2}C_K(S,\boldsymbol m)$ suffices.  No equality
case is asserted.

To prove this, suppose the stated quotient is infinite.  Its finite
Galois subextensions $K_i/K$ can then be chosen with degrees tending
to infinity.  They are unramified at every finite prime.  In each
finite Galois group the order of
$\operatorname{Frob}_{\mathfrak p}$ divides
$m_{\mathfrak p}$, and that order equals its residue degree.
Choose a subsequence whose signature ratios tend to
$\tau\in[0,1]$.  NT-H09 gives
$\varrho\ge e^{\pi\tau/2}C_K(S,\boldsymbol m)
\ge C_K(S,\boldsymbol m)$, contradicting the strict inequality.
In the totally real quotient every $K_i$ is totally real, so
$\tau=1$ and the second conclusion follows.  If $K=\mathbb Q$,
Minkowski's discriminant bound already makes the whole unramified
group trivial, so this omitted case satisfies the conclusion
unconditionally.  This completes the proof.

For $m_{\mathfrak p}=1$ this is the unrestricted closed-point
version of NT-H08.  If $S$ consists of all primes above rational
primes that split completely in $K$, then each such rational prime
contributes $n$ equal summands and NT-H08's constant is recovered.
For larger $m_{\mathfrak p}$ the imposed relations only bound the
orders of Frobenius elements.  The conclusion concerns the index
of their closed normal closure; it does not assert that this normal
subgroup is generated by finitely many elements as an ordinary
topological group.  This is a quantitative Galois-category
consequence of the TV explicit-formula inequality, with its precise
arithmetic input recorded, rather than a general anabelian
reconstruction statement.

\subsection{Symmetrization removes the infinity-type restriction for zero bounds}

\paragraph{NT-Z17. Uniform zero bounds for arbitrary unitary Hecke infinity type.}
Let $\eta$ be a primitive unitary Hecke character of a number field
$K$, $n=[K:\mathbb Q]$, which is not a pure norm character.
Write its standard completion as
\[
 \Lambda(s,\eta)=Q_\eta^{s/2}
         \prod_{j=1}^n\Gamma_{\mathbb R}(s+\mu_j)L(s,\eta),
 \qquad Q_\eta=\Delta_K\operatorname N\mathfrak f_\eta,
\]
up to a harmless nonzero constant.  Unitarity gives
$\operatorname{Re}\mu_j\ge0$, but the $\mu_j$ need not occur
in conjugate pairs.  For any real $T$, put
\[
 \mathscr C_\eta(T)=Q_\eta\prod_{j=1}^n(|\mu_j+iT|+3),
 \qquad
 R_\eta(T)=\log\log\bigl(\mathscr C_\eta(T)^{3/n}\bigr).
\]
Under GRH for $L(s,\eta)$, in particular under $M$,
\[
 \operatorname{ord}_{s=1/2+iT}L(s,\eta)
 \le\frac{\log\mathscr C_\eta(T)}{2R_\eta(T)}
  +O\left(\frac{\log\mathscr C_\eta(T)\log R_\eta(T)}
                 {R_\eta(T)^2}\right),
\tag{NT-Z17a}
\]
with an absolute implied constant.  Also, as
$\mathscr C_\eta(T)^{3/n}\to\infty$, there is a nontrivial zero
$1/2+i\gamma$ of this same $L$-function such that
\[
 |\gamma-T|\le\frac{\pi}{2R_\eta(T)}
                 +O\left(\frac{\log R_\eta(T)}{R_\eta(T)^2}\right),
\tag{NT-Z17b}
\]
again with an absolute implied constant.  The degree-normalized
conductor is retained, so these are uniform formulations even when
$K,\eta,T$, and $n$ vary within the stated class.

Here is the complete transfer proof.  Define the product and its
completion by
\[
 F_T(s)=L(s+iT,\eta)L(s-iT,\overline\eta),\qquad
 \mathcal F_T(s)=\Lambda(s+iT,\eta)
                         \Lambda(s-iT,\overline\eta).
\]
Classical Hecke theory makes both completed factors entire of order
one precisely because the pure norm exception has been excluded.
The gamma parameters of the product are
$\mu_j+iT$ and $\overline\mu_j-iT$, in conjugate pairs.
The factors $Q_\eta^{iT/2}$ and $Q_\eta^{-iT/2}$ cancel, so its
finite conductor is $Q_\eta^2$ and its degree is $2n$.
After grouping prime ideals by rational primes, the Euler parameters
of either shifted factor have modulus at most one.  Thus the product
has the parameter $\vartheta=0$ in CCM2015, \S4.  Conjugation and
the Hecke functional equation show that it is self-dual with root
number one.  GRH for $\eta$ implies GRH for its conjugate and for
both shifts, hence for the product.  It therefore satisfies all
the hypotheses of CCM2015 Theorems~5 and 7, including the actual
gamma-pairing condition.

Its analytic conductor in that source is exactly
\[
 C(\mathcal F_T)=Q_\eta^2
       \prod_{j=1}^n(|\mu_j+iT|+3)^2
       =\mathscr C_\eta(T)^2.
\]
Consequently its denominator is
\[
 \log\log\bigl(C(\mathcal F_T)^{3/(2n)}\bigr)
                  =R_\eta(T),
\]
with no alteration of degree dependence.
Conjugation gives
$\operatorname{ord}_{1/2-iT}L(s,\overline\eta)
=\operatorname{ord}_{1/2+iT}L(s,\eta)$.
The gamma factors are finite and nonzero at the points under
consideration, since their arguments have positive real part.
Thus the central order of $\mathcal F_T$ is twice the order on
the left of (NT-Z17a).  Apply CCM Theorem~5, equation (4.5): the
factor two in $\log C(\mathcal F_T)$ is cancelled by division
by this doubled order.  This gives (NT-Z17a), with a universal
error constant as stated.

By CCM Theorem~7 there is a zero $1/2+iu$ of $\mathcal F_T$
with the bound on $|u|$ appearing in (NT-Z17b).  If its first
factor vanishes, $L(s,\eta)$ has ordinate $T+u$.
If its second factor vanishes, conjugating that zero gives an
ordinate $T-u$ for $L(s,\eta)$.  Either case proves (NT-Z17b).
This argument also permits $u=0$, as it should.

For each fixed non-pure-norm $K,\eta$,
$\log\mathscr C_\eta(T)=n\log|T|+O_{K,\eta}(1)$ as
$|T|\to\infty$.  Substitution in the two bounds gives exactly
the multiplicity and centered covering-radius estimates of
NT-Z14 for this arbitrary infinity type.  Applying the covering
bound at the midpoint of a large gap first shows that its length
is less than one, and then gives
\[
 \gamma'-\gamma\le\frac{\pi}{\log\log\gamma}
    +O_{K,\eta}\left(
        \frac{\log\log\log\gamma}{(\log\log\gamma)^2}\right)
\]
for consecutive large positive ordinates.  If $\eta$ is a pure
norm character, its $L$-function is $\zeta_K(s+i\tau)$ for a
fixed real $\tau$.  Translate the finite-order trivial-character
case of NT-Z14 to obtain the same fixed-character asymptotic
conclusions.  We do not apply CCM verbatim to its two shifted
completed zeta factors: they have poles away from $0,1$ and would
violate that theorem's stated hypotheses.

This proves a zero-divisor transfer, not an individual argument
bound for $S_\eta(T)$ with unpaired gamma parameters.  The product
is the conjugate-dual pair of a rank-one arithmetic object and its
norm twist.  It provides a precise analytic use of that symmetry
without presupposing an algebraic motive for the infinite-order
character; compare root MT-H01 and the geometry programme's
self-dual determinant constructions.
When $L(s,\eta)=L(s+i\tau,\chi)$ for a primitive Dirichlet
character $\chi$, $A$ already supplies the needed zero-location
input by translation.

\paragraph{NT-Z18. At least half of the prime-conductor character family is nonvanishing.}
Let $q\to\infty$ run over primes and average over all $q-2$
nonprincipal characters modulo $q$.  Under $A$,
\[
 \frac1{q-2}\sum_{\chi\ne\chi_0}
             \operatorname{ord}_{s=1/2}L(s,\chi)
       \le\frac12+O\left(\frac1{\log q}\right),
\]
and therefore
\[
 \#\{\chi\ne\chi_0\pmod q:L(1/2,\chi)\ne0\}
       \ge\left(\frac12-O\left(\frac1{\log q}\right)\right)(q-2).
\tag{NT-Z18}
\]
For the full proof take
$f(x)=(\sin(2\pi x)/(2\pi x))^2$, with $f(0)=1$.
This is real, even, nonnegative, entire, and decays as
$O((1+|x|)^{-2})$ on the real line.  Its Fourier transform is
$\widehat f(u)=\tfrac12(1-|u|/2)_+$, supported on $[-2,2]$,
and $\int f=1/2$.  It meets Hughes--Rudnick2003 Definition~2.1,
author p.~4, including the analytic-strip condition in \S2.1.
Their Theorem~3.1, p.~5, permits this closed support and gives
average $W_f=1/2+O(1/\log q)$.
Under $A$ all the scaled ordinates in $W_f$ are real, so every
summand is nonnegative; the central zeros alone contribute their
multiplicity because $f(0)=1$.  This proves the mean-order bound.
The indicator of central vanishing is at most that multiplicity,
which proves (NT-Z18) after summing.  No parity assumption on
central multiplicity is used for this non-self-dual family.

The exact family and the positivity use of $A$ distinguish this
unitary-family corollary from the quadratic-family statement
NT-Z12.  In geometric language it concerns the central zero divisor
of all nontrivial rank-one objects of the prime cyclotomic cover;
it does not assert nonvanishing for each object.

\subsection{Classwise prime counts and the average conductor}

\paragraph{NT-C06. Explicit prime counts in a ray class.}
Let $K$ have degree $n$ and absolute discriminant $\Delta_K$.
Let $\mathfrak m=\mathfrak m_0\mathfrak m_\infty$ be a modulus,
$G=\operatorname{Cl}_{\mathfrak m}(K)$ its ray class group, and
$h=|G|$.  Write $H/K$ for its ray class field.  For
$\chi\in\widehat G$, let $\mathfrak f_\chi$ be the finite part
of the conductor of its primitive Hecke character.  Define
\[
 a_{\mathfrak m}=\log\Delta_K+
           \frac1h\sum_{\chi\in\widehat G}
                    \log\operatorname N\mathfrak f_\chi.
\]
The following unconditional relations will be used:
\[
 a_{\mathfrak m}=\frac{\log\Delta_H}{h},\qquad
 0\le a_{\mathfrak m}\le
 \log\Delta_K+\left(1-\frac1h\right)
                     \log\operatorname N\mathfrak m_0.
\tag{NT-ray-conductor}
\]
The lower inequality is only a convenient bound; the exact expression
retains the mean of the individual conductor logarithms.

For a class $c\in G$, put
\begin{align*}
 \psi_c^{\mathfrak m}(x)
 &=\sum_{\substack{\mathfrak p\nmid\mathfrak m_0,\ j\ge1\\
           (\operatorname N\mathfrak p)^j\le x,\ [\mathfrak p]^j=c}}
                  \log\operatorname N\mathfrak p,\\
 \pi_c^{\mathfrak m}(x)
 &=\#\{\mathfrak p\nmid\mathfrak m_0:
             \operatorname N\mathfrak p\le x,\ [\mathfrak p]=c\}.
\end{align*}
Let $e_{\mathfrak m}$ be the number of prime ideals dividing
$\mathfrak m_0$ which are unramified in $H/K$; in particular,
$e_{\mathfrak m}\le\omega_K(\mathfrak m_0)$.
Under RH for $\zeta_H$, and hence under $M$, one has
\begin{align*}
 \left|\psi_c^{\mathfrak m}(x)-\frac{x}{h}\right|
 &\le\sqrt x\left[
       \left(\frac{\log x}{2\pi}+2\right)a_{\mathfrak m}
        +\left(\frac{\log^2x}{8\pi}+2\right)n\right]
       +e_{\mathfrak m}\log x,
       &&x\ge1,                                      \tag{NT-C06a}\\
 \left|\pi_c^{\mathfrak m}(x)-\frac{\operatorname{Li}_2(x)}h\right|
 &\le\sqrt x\left[
       \left(\frac1{2\pi}+\frac3{\log x}\right)a_{\mathfrak m}
       +\left(\frac{\log x}{8\pi}+\frac1{4\pi}
                         +\frac6{\log x}\right)n\right]
       +e_{\mathfrak m},
       &&x\ge2.                                      \tag{NT-C06b}
\end{align*}
Here $\operatorname{Li}_2(x)=\int_2^xdt/\log t$.
Both estimates are uniform in the class $c$.
Their errors contain no multiplicative ray class number: its role
is retained in the main term and in the average conductor.
Replacing $a_{\mathfrak m}$ by the upper bound in
(NT-ray-conductor) gives a version using only
$K,\mathfrak m_0,h$.

We give the full deduction, with the exact classical inputs identified.
Class-field reciprocity and existence identify $G$ with
$\operatorname{Gal}(H/K)$ and the class of each
$\mathfrak p\nmid\mathfrak m_0$ with its Frobenius;
see Milne, \emph{Class Field Theory}, v4.03, V, Theorems~3.5--3.6,
printed p.~158/PDF p.~167.  The finite conductors of the associated
characters divide $\mathfrak m_0$, and the trivial character has
conductor one.  The conductor--discriminant formula gives
\[
 \Delta_H=\Delta_K^h
              \prod_{\chi\in\widehat G}
                      \operatorname N\mathfrak f_\chi.
\]
Its proof locator is Neukirch, \emph{Algebraic Number Theory},
VII, (11.7)--(11.9), printed pp.~533--534/PDF pp.~541--542;
(11.10), printed p.~535/PDF p.~543, identifies the degree-one
Artin and class-field conductors.  The complete supplied proof of
the conductor--discriminant formula on those pages was read on
the original images.  Taking logarithms proves the equality in
(NT-ray-conductor).  Bounding each of the $h-1$ nontrivial
conductors by $\mathfrak m_0$ proves its upper bound.

Every conjugacy class of $\operatorname{Gal}(H/K)$ is a singleton,
so its density is $1/h$.  Moreover $[H:\mathbb Q]=nh$.
Apply NT-C01/C02, the precise Greni\'e--Molteni2019
Theorem~1.1 and Corollary~1.2, to that singleton.  Dividing their
errors by $h$ replaces $\log\Delta_H$ by
$a_{\mathfrak m}$ and $[H:\mathbb Q]$ by $n$, giving exactly
the bracketed expressions above.

Those source counts include every prime unramified in $H/K$.
All primes outside $\mathfrak m_0$ are unramified, but a prime
inside $\mathfrak m_0$ may be unramified as well.  Each such prime
contributes at most
\[
 \sum_{\substack{j\ge1\\(\operatorname N\mathfrak p)^j\le x}}
          \log\operatorname N\mathfrak p\le\log x
\]
to a prime-power count, and at most one to an ordinary prime
count.  Removing these $e_{\mathfrak m}$ possible primes and
using the triangle inequality proves (NT-C06a)--(NT-C06b), including
degenerate moduli.  This completes the proof.  For example, the
prime $2$ is unramified in the ray class field $\mathbb Q$ for
modulus $2\infty$ and is correctly removed by this boundary term.

The input is Dedekind RH for this one field $H$; no Artin
holomorphy conjecture is needed.  When $H/\mathbb Q$ is abelian,
$A$ suffices by Dirichlet factorization.  In particular, $A$
supplies every base-field-$\mathbb Q$ ray-class specialization.
For general $K$, abelianity of $K/\mathbb Q$ alone does not assert
abelianity of $H/\mathbb Q$ and is not substituted for this input.
Geometrically, the count takes place in the finite abelian
quotient of $\pi_1^{\text{ét}}(\operatorname{Spec}
\mathcal O_K[1/\mathfrak m_0])$ specified by the modulus and the
real-place conditions.  The quantity $a_{\mathfrak m}$ is the
average Artin conductor cost of its rank-one objects.

\paragraph{NT-C07. Ideal classes and degree-one prime representatives.}
Let $h_K=|\operatorname{Cl}(K)|$ and assume RH for
$\zeta_{\mathbb Q}$ and the Hilbert class-field Dedekind function.  Then every ordinary ideal class
contains a prime ideal of residue degree one over $\mathbb Q$
whose norm is at most
\[
 \left[h_K\bigl(4\log\Delta_K+2.5n\bigr)+5\right]^2.
\tag{NT-C07}
\]
For the narrow ideal class group the same conclusion holds with
$h_K^+$ in place of $h_K$, assuming RH for the corresponding
narrow class field.  $M$ supplies either hypothesis.
For either group, the classwise estimates (NT-C06a)--(NT-C06b)
hold with $a_{\mathfrak m}=\log\Delta_K$ and
$e_{\mathfrak m}=0$.

Indeed, use the modulus with finite part one and with empty real
part for the ordinary group, or all real places for the narrow
group.  Its class field $H/K$ is unramified at all finite primes;
consequently $\Delta_H=\Delta_K^{[H:K]}$.
The prime-count assertions follow at once from NT-C06.
For $K\ne\mathbb Q$ one has $H\ne\mathbb Q$, so
Bach--Sorenson1996 Theorem~5.1, printed p.~1729/PDF p.~13,
applies to each singleton of $\operatorname{Gal}(H/K)$.
Its degree-one prime has norm at most
$(4\log\Delta_H+2.5[H:\mathbb Q]+5)^2$.
Substitute $[H:K]=h_K$ or $h_K^+$ and the discriminant equality
to obtain (NT-C07).  The source requires RH for $\zeta_H$
and $\zeta_{\mathbb Q}$; $M$ provides both, and $A$ already
provides the latter.  If $K=\mathbb Q$, both groups are trivial
and the prime $2$ directly proves the displayed bound, without
applying the source theorem to its excluded field.
This proves the assertion in all cases.

The ordinary class group is
$\operatorname{Pic}(\operatorname{Spec}\mathcal O_K)$; the narrow
version adds the real sign conditions.  These are prime
representatives of each individual class, distinct from a small
set of primes generating the whole group.  No arbitrary extra
forbidden-prime set is included in (NT-C07); that additional
condition can instead be tested by the explicit count in NT-C06.

\paragraph{NT-C08. Certified ray-class Mertens constants and Euler products.}
Retain the notation and exact RH hypothesis of NT-C06, and let
$\delta=1/h$.  For each ray class $c$ there are real constants
$B_c$ and $A_c>0$ such that, for every $x\ge2$,
\begin{align*}
 \left|\sum_{\substack{\mathfrak p\nmid\mathfrak m_0\\
          [\mathfrak p]=c,\ \operatorname N\mathfrak p\le x}}
        \frac1{\operatorname N\mathfrak p}
        -\delta\log\log x-B_c\right|
       &\le U_{\mathfrak m}(x),                       \tag{NT-C08a}\\
 \left|\log\prod_{\substack{\mathfrak p\nmid\mathfrak m_0\\
          [\mathfrak p]=c,\ \operatorname N\mathfrak p\le x}}
       (1-(\operatorname N\mathfrak p)^{-1})^{-1}
       -\delta\log\log x-\log A_c\right|
       &\le U_{\mathfrak m}(x)+\frac{n}{2\lfloor x\rfloor},
                                                        \tag{NT-C08b}
\end{align*}
where the explicit error is
\[
 U_{\mathfrak m}(x)=\frac1{\sqrt x}
 \left[
   \left(\frac3{2\pi}+\frac9{\log x}\right)a_{\mathfrak m}
   +\left(\frac{3\log x}{8\pi}+\frac5{4\pi}
                    +\frac{18}{\log x}\right)n\right]
       +\frac{2e_{\mathfrak m}}x.
\tag{NT-C08-error}
\]
In particular the product divided by $(\log x)^{1/h}$ tends
to $A_c$.  The two finite expressions in (NT-C08a)--(NT-C08b),
with these errors, give certified intervals for $B_c$ and
$\log A_c$ once the relevant primes have been enumerated.
No uncomputed numerical value of either constant is asserted.

Here is the complete proof.  Abbreviate
$\pi_c^{\mathfrak m}(t)=\delta\operatorname{Li}_2(t)+R_c(t)$,
and let $S_c(x)$ be the reciprocal-prime sum in (NT-C08a).
Partial summation, including possible primes of norm $2$, gives
\[
 S_c(x)=\frac{\pi_c^{\mathfrak m}(x)}x+
                  \int_2^x\frac{\pi_c^{\mathfrak m}(t)}{t^2}\,dt.
\]
The lower boundary uses $\pi_c^{\mathfrak m}(2^-)=0$;
at $x=2$ the displayed formula is still valid.
Since $\operatorname{Li}_2(2)=0$, integration by parts yields
\[
 \frac{\operatorname{Li}_2(x)}x+
          \int_2^x\frac{\operatorname{Li}_2(t)}{t^2}\,dt
               =\log\log x-\log\log2.
\]
The estimate in NT-C06 makes $R_c(t)/t^2$ absolutely integrable
on $[2,\infty)$.  Define
\[
 B_c=-\delta\log\log2+
                       \int_2^\infty\frac{R_c(t)}{t^2}\,dt.
\]
These identities give the exact error formula
\[
 S_c(x)-\delta\log\log x-B_c
          =\frac{R_c(x)}x-
                        \int_x^\infty\frac{R_c(t)}{t^2}\,dt.
\tag{NT-C08-tail}
\]
Apply NT-C06b to both terms.  The elementary integrals needed are
\[
 \int_x^\infty t^{-3/2}\,dt=\frac2{\sqrt x},\qquad
 \int_x^\infty t^{-3/2}\log t\,dt
                      =\frac{2\log x+4}{\sqrt x},\qquad
 \int_x^\infty\frac{t^{-3/2}}{\log t}\,dt
                      \le\frac2{\sqrt x\log x}.
\]
For the $a_{\mathfrak m}$ terms the point and integral bounds
sum to $3/(2\pi)+9/\log x$ after division by $\sqrt x$.
For the $n$ terms they sum to
$3\log x/(8\pi)+5/(4\pi)+18/\log x$.
The boundary-prime contribution is
$e_{\mathfrak m}/x+\int_x^\infty e_{\mathfrak m}t^{-2}dt
=2e_{\mathfrak m}/x$.  This proves (NT-C08a) with exactly
the error displayed.

To pass to the product, define the absolutely convergent series
\[
 D_c=\sum_{\substack{\mathfrak p\nmid\mathfrak m_0\\
                         [\mathfrak p]=c}}
                 \sum_{k\ge2}\frac1{k(\operatorname N\mathfrak p)^k},
 \qquad A_c=\exp(B_c+D_c).
\]
For $q\ge2$,
$\sum_{k\ge2}1/(kq^k)\le1/(2q(q-1))$.
There are at most $n$ prime ideals of $K$ of any given
integer norm: if that norm is a rational prime power $p^f$,
there are at most $n/f\le n$ primes of that residue degree,
by the degree formula over $p$, and otherwise there are none.
Consequently
\[
 0\le\sum_{\substack{\mathfrak p\nmid\mathfrak m_0\\
              [\mathfrak p]=c,\ \operatorname N\mathfrak p>x}}
            \sum_{k\ge2}\frac1{k(\operatorname N\mathfrak p)^k}
 \le\frac n2\sum_{j>\lfloor x\rfloor}\frac1{j(j-1)}
 =\frac{n}{2\lfloor x\rfloor}.
\]
The logarithm of the finite Euler product is $S_c(x)$ plus
the corresponding truncated $k\ge2$ series.  Subtract
$\delta\log\log x+B_c+D_c$, use (NT-C08a) and this last
tail estimate, and obtain (NT-C08b).  Both errors tend to zero;
exponentiation proves the positive limiting constant and the
claimed interval certification.  This completes the proof.

The same deductions apply to ordinary and narrow ideal classes
with $a_{\mathfrak m}=\log\Delta_K$ and no boundary term.
For $K=\mathbb Q$ they are consequences of $A$ for every ray
class modulus.  The geometric object is the reciprocal-norm
measure on the closed points with a prescribed image in the
abelian fundamental-group quotient; its regularized Euler
product has a quantitatively controlled real constant.
This extends the finite prime counts to a convergent
regularization, without identifying that constant with a
particular motivic period.

\subsection{A conductor-uniform critical-line bound}

\paragraph{NT-L13. Uniform Lindel\"of bounds for non-pure-norm unitary Hecke characters.}
Let $\eta$ be a primitive unitary Hecke character of $K$ which is
not a pure norm character, and let $n=[K:\mathbb Q]$.
Use the completion and the sharp conductor already defined in
NT-Z17:
\[
 \mathscr C_\eta(T)=Q_\eta
                  \prod_{j=1}^n(|\mu_j+iT|+3),
 \qquad y=\log\log\mathscr C_\eta(T).
\]
Assume GRH for $L(s,\eta)$ and RH for $\zeta_{\mathbb Q}$.
Thus $M$ supplies both inputs; $A$ suffices for a norm shift
of a nonprincipal primitive Dirichlet character.
If $y\ge10$, then
\[
 |L(1/2+iT,\eta)|\le
  \exp\left[
    \left(\frac3{8y}+\frac{23n}{25y^2}\right)
                      \log\mathscr C_\eta(T)\right].
\tag{NT-L13}
\]
All constants are explicit and independent of $K,\eta,T,n$.
In particular, for every $\epsilon>0$ one has
$|L(1/2+iT,\eta)|\le\mathscr C_\eta(T)^\epsilon$ whenever
\[
 y\ge\max\left\{10,\frac3{4\epsilon},
                         \sqrt{\frac{46n}{25\epsilon}}\right\}.
\tag{NT-L13-threshold}
\]
This gives the conductor-uniform Lindel\"of statement for fixed
degree.  If the degree varies, the same conclusion in the form
$|L|\le\mathscr C^{o(1)}$ follows along any family for which
$y\to\infty$ and $n=o(y^2)$.  The dependence on degree is thus
specified rather than absorbed in an implied constant.

We verify the full specialization and the threshold deduction.
The classical input is Chandee2009, \emph{Explicit upper bounds
for $L$-functions on the critical line}, Corollary~1.1,
author p.~3, whose axioms and conductor are on p.~2,
equations (2)--(6).  Its proof is on p.~7, equations
(14)--(15); the underlying logarithmic inequality is proved
in Theorem~2.1, pp.~4--6.
Apply it to $s\mapsto L(s+iT,\eta)$.
The gamma parameters are $\mu_j+iT$ and have nonnegative
real part; this source does not impose the gamma-pairing
condition encountered in CCM2015.  Standard primitive Hecke
continuation and the functional equation give its analytic
axioms, and non-pure-norm unitary Hecke functions have no
completed pole or zero on the boundary lines.  The latter
nonvanishing is also part of the adopted GRH zero-location
input; its unconditional classical form is recorded in
Ramakrishnan--Valenza, Theorem~7-28, printed p.~289.

Grouping the prime-ideal Euler factors over each rational
prime produces $n$ parameters of modulus at most one,
padding with zero parameters at ramified primes when needed.
Indeed a factor $1-\eta(\mathfrak p)p^{-fs}$ factors into
$f$ linear factors in $p^{-s}$ whose roots have modulus one;
the degree formula gives total degree at most $n$.
The height shift multiplies those roots by $p^{-iT}$ and
preserves their moduli.  This proves the Ramanujan hypothesis
used in the corollary.  It also proves that the logarithmic
Euler coefficient at $p^k$ has modulus at most $n\log p$.
The source's proof of its explicit prime-power estimate uses
RH for $\zeta_{\mathbb Q}$, which is why that input has been
stated separately above.

In the source's normalization the conductor is
\[
 C_{\mathrm{Ch}}=
 \frac{Q_\eta}{\pi^n}
        \prod_{j=1}^n\left|\frac14+
                              \frac{\mu_j+iT}{2}\right|
 =\frac{Q_\eta}{(2\pi)^n}
        \prod_{j=1}^n|1/2+\mu_j+iT|
 \le\mathscr C_\eta(T).
\]
It is legitimate to use this larger conductor even if only
its double logarithm exceeds ten.  To see this directly in
the cited proof, retain the free parameter $x$ in equation
(14), replace $\log C_{\mathrm{Ch}}$ by
$\log\mathscr C_\eta(T)$ in its term with positive
coefficient $3/(4\log x)$, and then choose
$x=(\log\mathscr C_\eta(T))^2$.
The estimate (15) and the remaining terms depend only on
$x$ and $n$, and $\log x=2y\ge20$ meets their threshold.
The same collection of constants that proves Corollary~1.1
therefore gives exactly (NT-L13).  This explains the conductor
normalization and does not silently replace it by an
asymptotically equivalent one when the degree varies.
If the central value is zero, the nonlogarithmic bound is
immediate; otherwise the logarithmic argument applies.

Finally the two conditions
$y\ge3/(4\epsilon)$ and
$y^2\ge46n/(25\epsilon)$ make the two terms
$3/(8y)$ and $23n/(25y^2)$ at most $\epsilon/2$ each.
Exponentiation proves (NT-L13-threshold).
The varying-degree assertion follows from the same two
terms tending to zero.  This completes the deduction.

Pure norm characters have shifted completed zeta poles and
are not inserted into these pole-free axioms; their fixed-function
Lindel\"of consequences remain in NT-L03.  Geometrically,
(NT-L13) bounds the critical value of a rank-one unitary
arithmetic object uniformly in its finite and archimedean
conductors.  No algebraicity of an infinite-order infinity
type, negative-moment estimate, or nonzero lower bound at the
critical point is required or inferred.

\paragraph{NT-L14. Explicit two-sided bounds on the one-line.}
Let $\eta$ be a primitive unitary non-pure-norm Hecke character
over a degree-$n$ field, and set
\[
 C_1(\eta,T)=\frac{Q_\eta}{(2\pi)^n}
                   \prod_{j=1}^n|1+\mu_j+iT|.
\]
Here the exact normalization is retained.
For $\log C\ge23n$ put $Y=\log\log C-\log(2n)>0$ and define
\begin{align*}
 k_n&=2.31+\frac{22.59}{n}(e^{0.31n}-1-0.31n),\\
 j_n&=2+\frac{4.18}{n}(e^{0.69n}-1-0.69n),\\
 j_n'&=9+\frac{16.74}{n}(e^{0.69n}-1-0.69n),\\
 U_n(C)&=(2e^\gamma)^n
       \left(Y^n+\frac n2Y^{n-1}+\frac{nk_n}4Y^{n-2}\right),\\
 V_n(C)&=\left(\frac{12e^\gamma}{\pi^2}\right)^n
       \left(Y^n+\frac n2Y^{n-1}+\frac{nj_n}4Y^{n-2}
                         +\frac{n^2j_n'Y^n}{\log C}\right).
\end{align*}
Under GRH for $L(s,\eta)$ and RH for $\zeta_{\mathbb Q}$,
\[
  V_n(C_1(\eta,T))^{-1}
        \le |L(1+iT,\eta)|\le U_n(C_1(\eta,T))
        \qquad\text{if }\log C_1(\eta,T)\ge23n.
\tag{NT-L14}
\]
Thus $M$ supplies the hypotheses, while $A$ supplies them for
norm shifts of nonprincipal primitive Dirichlet characters.
For fixed $n$ the leading upper order is
$(2e^\gamma\log\log C)^n$, and the reciprocal leading order
is $(12e^\gamma\log\log C/\pi^2)^n$.
These are two-sided inequalities, not a claim that every value
has either extremal order.

This is Lumley2018, \emph{Explicit bounds for $L$-functions
on the edge of the critical strip}, Theorems~2.1--2.2,
author pp.~3--4, and Corollary~2.3, p.~5.
The definitions and axioms are in \S1.1, pp.~2--3;
the proofs are \S4.1, pp.~11--13, and \S4.2, pp.~14--16.
The exact constants (2.1)--(2.3) and both theorem formulas
were inspected on the original p.~4 image.
The rational-prime Euler-factor and analytic-axiom verification
is exactly the one given in NT-L13, with parameters shifted
by $iT$.  The source allows complex infinity parameters and
requires an entire completion, which explains the non-pure-norm
condition.  Its Lemmas~3.6--3.7, p.~10, explicitly use RH for
$\zeta_{\mathbb Q}$; that input is retained here.
Its quantity $J_1(n)$ is at most $j_n$, while $J_2(n)=j_n'$.
Since $Y>0$, replacing $J_1(n)$ by its displayed upper bound
preserves the reciprocal inequality.  This verifies every
specialization and normalization in (NT-L14).
The numerical constants remain the cited classical theorem;
no new numerical optimization is being claimed.

\paragraph{NT-H11. Relative class-number--regulator bounds over any base field.}
Let $E/K$ be a finite abelian extension with group $G$, and
write $n=[K:\mathbb Q]$.
For every $\chi\in\widehat G\setminus\{1\}$, let
$L(s,\chi)$ mean the primitive Hecke function attached to that
Artin character, including its correct ramified local factors.
Let $\mathfrak f_\chi$ be its finite conductor and let
$a_\chi$ count the real places of $K$ where its sign is odd.
Its conductor for NT-L14 at height zero is explicitly
\[
 C_\chi=
 \frac{\Delta_K\operatorname N\mathfrak f_\chi}
                  {(2\pi)^n}\,2^{a_\chi+r_2(K)}.
\tag{NT-relative-conductor}
\]
For a number field $F$, let $h_F$ be the ordinary ideal class
number, $R_F$ its full unit regulator, $w_F$ its number of roots
of unity, and $r_1(F),r_2(F)$ its signature.  Put
\[
 D(E/K)=\frac{w_E}{w_K}
       \sqrt{\frac{\Delta_E}{\Delta_K}}\,
       2^{r_1(K)-r_1(E)}(2\pi)^{r_2(K)-r_2(E)}.
\]
Assume RH for $\zeta_E$ and $\zeta_{\mathbb Q}$; in particular,
$M$ suffices.  If every nontrivial character satisfies
$\log C_\chi\ge23n$, then
\[
 D(E/K)\prod_{\chi\ne1}V_n(C_\chi)^{-1}
 \le\frac{h_ER_E}{h_KR_K}
 \le D(E/K)\prod_{\chi\ne1}U_n(C_\chi).
\tag{NT-H11}
\]
If some conductors fall below this threshold, let
$\mathcal S$ be precisely that finite set of characters and
put $F_{\mathrm{small}}=\prod_{\chi\in\mathcal S}|L(1,\chi)|$.
The valid bounds are then (NT-H11) with the products restricted
to $\chi\notin\mathcal S\cup\{1\}$ and both endpoints
multiplied by $F_{\mathrm{small}}$.
Thus no small-conductor value has been silently set to one.
A bound on $h_E/h_K$ alone additionally requires information
on $R_E/R_K$.

Here is the complete proof, including the ramified factors.
At a prime $\mathfrak p$ of $K$ put $q=\operatorname N\mathfrak p$.
Let $I\subset G$ be its inertia group, of order $e$, and let
$g\in G/I$ be Frobenius of order $f$.
There are $|G|/(ef)$ primes of $E$ above $\mathfrak p$,
all of norm $q^f$; the local Dedekind factor of $E$ is therefore
$(1-q^{-fs})^{-|G|/(ef)}$.
Only the characters trivial on $I$ contribute a local Hecke
factor, namely $(1-\chi(g)q^{-s})^{-1}$.
On restricting the characters of $G/I$ to its cyclic subgroup
$\langle g\rangle$, each of the $f$ roots of unity occurs
$|G|/(ef)$ times.  Consequently the product of these local
Hecke factors is the same Dedekind factor.
This proves on the absolutely convergent half-plane, and hence
by continuation,
\[
 \zeta_E(s)=\zeta_K(s)\prod_{\chi\ne1}L(s,\chi).
\tag{NT-relative-factorization}
\]
This uses the classical class-field identification of finite-order
Hecke characters with the characters of $G$, as in Milne V,
Theorems~3.5--3.6; the local factor computation above proves
the actual Euler-product identity used here.

All the nontrivial Hecke factors are entire, and the only pole
of $\zeta_K$ is at $1$.  A nontrivial zero off the critical
line of any factor would therefore be a nontrivial zero of
$\zeta_E$.  Thus RH for $\zeta_E$ supplies the Hecke GRH
hypotheses required by NT-L14.  The remaining RH for
$\zeta_{\mathbb Q}$ has been stated explicitly.
At each real place the gamma parameter of $\chi$ is $0$ or $1$;
at a complex place the two real-gamma parameters are $0,1$.
Substitution into the definition of $C_1$ proves
(NT-relative-conductor), with no asymptotic normalization.

Let $\kappa_F=\operatorname{Res}_{s=1}\zeta_F(s)>0$.
Taking residues in (NT-relative-factorization) and absolute
values gives
\[
 \frac{\kappa_E}{\kappa_K}
                  =\prod_{\chi\ne1}|L(1,\chi)|.
\]
The analytic class-number formula in Milne VI,
Corollary~2.12 and Example~2.13, printed p.~191/PDF p.~200,
is
\[
 \kappa_F=\frac{2^{r_1(F)}(2\pi)^{r_2(F)}h_FR_F}
                         {w_F\sqrt{\Delta_F}}.
\]
Its ratio therefore yields the exact identity
\[
 \frac{h_ER_E}{h_KR_K}
             =D(E/K)\prod_{\chi\ne1}|L(1,\chi)|.
\]
Apply NT-L14 to each factor above its threshold and leave the
remaining factors equal to their actual values.  This proves
both (NT-H11) and the small-conductor version.
Empty products give the correct identity when $E=K$.
The proof is complete.

For $E/\mathbb Q$ abelian, $A$ supplies RH for $\zeta_E$,
recovering and sharpening the characterwise framework of NT-H04.
For a general abelian extension $E/K$, the needed Dedekind
input is genuinely that of $E$, not just that of $K$.
The geometric comparison is between the Picard groups of the
two arithmetic curves and the rank-one character decomposition
of their finite abelian cover.  Its regulator and real-place
factors are integral parts of the comparison, not omitted
normalization terms.

\subsection*{NT-L15. A mean-square criterion for simple Hecke zeros and reciprocal derivatives}
Let $\eta$ be a fixed primitive unitary Hecke character of a number field $K$,
including a pure norm character. Define coefficients $b_\eta(\mathfrak a)$ by
\[
 L(s,\eta)^{-1}=\sum_{\mathfrak a\ne0}
       b_\eta(\mathfrak a)(N\mathfrak a)^{-s}
                    \qquad(\Re s>1),
 \qquad
 \mathcal M_\eta(x)=\sum_{N\mathfrak a\le x}b_\eta(\mathfrak a).
\]
Thus $b_\eta(\mathfrak a)=\mu_K(\mathfrak a)\eta(\mathfrak a)$
on ideals prime to the conductor and is zero on the other ideals.
Assume the additional weak mean-square hypothesis
\[
 B_\eta:=\limsup_{X\to\infty}\frac1{\log X}
       \int_1^X\left|\frac{\mathcal M_\eta(x)}x\right|^2dx<\infty.
                                                        \tag{NT-L15a}
\]
Then GRH holds for $L(s,\eta)$, every nontrivial zero is simple, and
\[
 \sum_{\rho=1/2+i\gamma}\frac1{|\rho L'(\rho,\eta)|^2}
                          \le B_\eta.                   \tag{NT-L15b}
\]
In particular $B_\eta>0$ and
\[
 |L'(\rho,\eta)|^{-1}\le\sqrt{B_\eta}\,|\rho|,
 \qquad |L'(\rho,\eta)|^{-1}=o(|\rho|),
                                                        \tag{NT-L15c}
\]
where the little-$o$ is along zeros of unbounded height. Also
\[
 \sum_{|\gamma|\le T}|L'(1/2+i\gamma,\eta)|^{-2}=o(T^2).
                                                        \tag{NT-L15d}
\]
No use of $A$ or $M$ is required once (NT-L15a) is assumed. Neither $A$ nor
$M$ is asserted to imply this additional mean-square hypothesis.
For $K=\mathbb Q$, $\eta=1$, the classical source is Titchmarsh,
\emph{The Theory of the Riemann Zeta-function}, 2nd ed.,
\S14.29: hypothesis (14.29.1) and Theorem14.29(A), pp.376--377;
reciprocal derivative estimate (14.29.3), p.377;
Theorem14.29(B)/(14.29.4), p.377, with proof ending p.380.
The following proof gives the Hecke formulation and its precise limsup constant.

\begin{proof}
Put $F(y)=e^{-y/2}\mathcal M_\eta(e^y)$ for $y\ge0$ and
$A(Y)=\int_0^Y|F(y)|^2dy$. Changing variables $x=e^y$ gives
$B_\eta=\limsup_{Y\to\infty}A(Y)/Y$.
For every $\epsilon>0$ there is $C_\epsilon$ such that
\[
 A(Y)\le (B_\eta+\epsilon)Y+C_\epsilon\qquad(Y\ge0).
                                                        \tag{NT-L15e}
\]
In particular the Laplace integral
\[
 H(w)=\int_0^\infty F(y)e^{-wy}\,dy
\]
is holomorphic for $\Re w>0$. To check this, integrate $|F|^2$ against
$e^{-\sigma y}$ by parts using (NT-L15e), then apply Cauchy--Schwarz to
$F(y)e^{-wy}$. The same argument with any fixed factor $y^j$ proves
uniform convergence of every derivative on each compact subset of that
half-plane. For $\Re w>1/2$, absolute convergence of the inverse Euler
series and partial summation give
\[
 H(w)=\int_1^\infty\mathcal M_\eta(x)x^{-(w+1/2)-1}\,dx
      =\frac1{(w+1/2)L(w+1/2,\eta)}.                    \tag{NT-L15f}
\]
The meromorphic identity theorem extends (NT-L15f) throughout $\Re w>0$.
Since $H$ is holomorphic there and $w+1/2\ne0$, $L(s,\eta)$ has no zeros
with $\Re s>1/2$. The conjugate character has coefficients
$b_{\bar\eta}(\mathfrak a)=\overline{b_\eta(\mathfrak a)}$, hence the
same mean-square hypothesis and the same zero-free half-plane.
The Hecke functional equation now excludes nontrivial zeros of
$L(s,\eta)$ with $\Re s<1/2$. This proves GRH, including non-real
characters; possible pure-norm poles only give zeros of $H$ and cause
no difficulty.

For $\sigma>0$ introduce the positive inner product
\[
 \langle f,g\rangle_\sigma
       =\sigma\int_0^\infty f(y)\overline{g(y)}e^{-\sigma y}dy.
\]
An integration by parts and (NT-L15e) show that
\[
 \|F\|_\sigma^2
   =\sigma^2\int_0^\infty A(y)e^{-\sigma y}dy,
 \qquad \limsup_{\sigma\downarrow0}\|F\|_\sigma^2\le B_\eta.
                                                        \tag{NT-L15g}
\]
Each $e_t(y)=e^{ity}$ has norm one. Consequently
\[
 |\sigma H(\sigma+it)|
       =|\langle F,e_t\rangle_\sigma|
       \le\|F\|_\sigma.
                                                        \tag{NT-L15h}
\]
At a zero $\rho=1/2+i\gamma$, (NT-L15f) has a pole of the same order
as that zero, since $\rho\ne0$. The bound (NT-L15h) as
$\sigma\downarrow0$ rules out order greater than one. Thus every zero
is simple and
\[
 \lim_{\sigma\downarrow0}\langle F,e_\gamma\rangle_\sigma
       =a_\gamma:=\frac1{\rho L'(\rho,\eta)}.
                                                        \tag{NT-L15i}
\]
For two distinct ordinates,
\[
 \langle e_\gamma,e_\nu\rangle_\sigma
       =\frac\sigma{\sigma-i(\gamma-\nu)}\longrightarrow0;
\]
the same inner product is one when $\gamma=\nu$.
Take an arbitrary finite set $S$ of zero ordinates and put
$P_S=\sum_{\gamma\in S}a_\gamma e_\gamma$.
Expanding the nonnegative number $\|F-P_S\|_\sigma^2$, using
(NT-L15g)--(NT-L15i), and taking the upper limit gives
\[
 0\le\limsup_{\sigma\downarrow0}\|F-P_S\|_\sigma^2
       \le B_\eta-\sum_{\gamma\in S}|a_\gamma|^2.
\]
The supremum over finite $S$ proves (NT-L15b). Nontrivial Hecke zeros
exist, by the classical zero-count asymptotic, so $B_\eta>0$.
Every summand is at most $B_\eta$, proving the first assertion of
(NT-L15c). Square summability makes $a_\gamma\to0$ along unbounded
ordinates and proves the second assertion. Finally, for $T\ge1$,
\[
 T^{-2}\sum_{|\gamma|\le T}|L'(\rho,\eta)|^{-2}
   =\sum_\gamma {\bf1}_{|\gamma|\le T}
           \frac{1/4+\gamma^2}{T^2}|a_\gamma|^2.
\]
Each factor before $|a_\gamma|^2$ tends to zero and is at most $5/4$.
Dominated convergence for the summable series proves (NT-L15d).
\end{proof}

The finite Bessel argument proves an inequality. It does not assert an
exact mean-square limit, completeness of an infinite Fourier expansion
of $F$, or a limiting distribution for $F$. Such assertions require
additional approximation estimates.
Geometrically, the coefficients are the inverse Euler weights of the
rank-one arithmetic object attached to $\eta$; squarefree ideals are
squarefree effective divisors on $\operatorname{Spec}\mathcal O_K$.
The new hypothesis is analytic control of these weighted divisors and
is not a consequence of a finite-field purity theorem established here.

\subsection*{NT-Z19. Weak mean square supplies the canonical arithmetic basis gate}
Use the canonical completed, pole-removed $\Xi_\eta$, inverse Mellin
function $\psi_\eta$, and closed quotient
\[
 Q_\eta=\mathcal E/\overline{\psi_\eta*\mathcal E},\qquad
 \mathcal E=\{f\in C^\infty(\mathbb R):
      \sup_x e^{N|x|}|f^{(j)}(x)|<\infty\text{ for all }N,j\ge0\}
\]
of MT-M12/MT-M16. Under (NT-L15a), all zeros are simple and
\[
 \overline{\psi_\eta*\mathcal E}
   =\{f:\mathcal Mf(\rho)=0\text{ for every nontrivial zero }\rho\},
 \qquad Q_\eta\simeq s(Z).
                                                        \tag{NT-Z19a}
\]
Here $\mathcal Mf(s)=\int_\mathbb Rf(x)e^{sx}dx$,
and $s(Z)$ is the space of sequences decreasing faster than every
power of $1+|\operatorname{Im}\rho|$. The quotient has an absolute
Schauder basis with these Mellin evaluations as coefficient functionals.
Convolution by every $h\in\mathcal E$ is nuclear. In the diagonal
Hilbert realization $\ell^2(Z)$ it is trace class, with Hilbert trace
\[
 \operatorname{tr}_{\ell^2(Z)}T_h=\sum_\rho\mathcal Mh(\rho).
\]
Writing $h^\#(x)=e^{-x}\overline{h(-x)}$, the operator
$T_{h*h^\#}$ is positive and
\[
 \operatorname{tr}_{\ell^2(Z)}T_{h*h^\#}
                     =\sum_\rho|\mathcal Mh(\rho)|^2\ge0.
                                                        \tag{NT-Z19b}
\]
\begin{proof}
NT-L15 proves GRH, simplicity, and
$|L'(\rho,\eta)|\ge B_\eta^{-1/2}|\rho|^{-1}$.
Since $|\rho|\le1+|\operatorname{Im}\rho|$, this is precisely the
polynomial derivative lower bound of MT-M12, or condition (1) of
MT-M16 with every multiplicity equal to one. The entire proofs of
those results give (NT-Z19a) and the absolute basis. MT-M13/MT-M17
then give nuclearity and the trace-class Hilbert realization. For
completeness, on each coordinate the convolution eigenvalue is
$\mathcal Mh(\rho)$, and direct substitution gives
$\mathcal M h^\#(s)=\overline{\mathcal Mh(1-\bar s)}$.
GRH implies $1-\bar\rho=\rho$, so the product eigenvalue is
$|\mathcal Mh(\rho)|^2$. Rapid Mellin decay and NT-Z10's local zero
count make both displayed trace series absolutely convergent. This
proves positivity and (NT-Z19b).
\end{proof}
This is a conditional analytic realization of the zero divisor with an
explicit arithmetic gate. It does not identify the quotient with an
independently constructed motivic or étale cohomology group.
For the Riemann zeta function the sole additional arithmetic condition
is the classical weak Mertens conjecture; $A$ is not needed separately.
